\documentclass[11pt,letterpaper]{article}
\usepackage[T1]{fontenc}
\usepackage{lmodern,microtype}
\usepackage[margin=1.02in,headheight=14pt]{geometry}
\usepackage{amsmath,amssymb,amsthm,mathtools,booktabs,needspace}
\usepackage{xcolor,fancyhdr}
\definecolor{teal}{HTML}{175463}
\usepackage[colorlinks=true,linkcolor=teal,citecolor=teal,urlcolor=teal]{hyperref}
\usepackage{bookmark}
\hypersetup{pdftitle={Explicit zero locations in the Newman heat flow},pdfauthor={Drafted for Edward Baker}}
\pagestyle{fancy}\fancyhf{}
\fancyhead[L]{\small Explicit zero locations in the Newman heat flow}
\fancyhead[R]{\small Working draft}
\fancyfoot[C]{\thepage}
\setlength{\headsep}{16pt}
\setlength{\emergencystretch}{1.5em}
\raggedbottom
\numberwithin{equation}{section}
\setcounter{tocdepth}{1}
\newtheorem{theorem}{Theorem}[section]
\newtheorem{proposition}[theorem]{Proposition}
\newtheorem{lemma}[theorem]{Lemma}
\newcommand{\E}{\mathbb E}
\newcommand{\ii}{\mathrm i}
\newcommand{\dd}{\,\mathrm d}
\newcommand{\Real}{\operatorname{Re}}
\newcommand{\Imag}{\operatorname{Im}}
\newcommand{\Var}{\operatorname{Var}}
\newcommand{\Res}{\operatorname{Res}}
\begin{document}
\thispagestyle{plain}
\begin{center}
{\small\scshape Analytical working draft}\\[8pt]
{\LARGE\bfseries Explicit zero locations in the Newman heat flow}\\[5pt]
{\large Counted positions, high ranks, and backward gap curvature}\\[9pt]
{\large Drafted for Edward Baker}\\[3pt]
20 November 2026
\end{center}
\begin{abstract}
We derive estimates for consecutive zeros of the Newman heat flow from the theta source at one time. As the time \(T\) grows, a single inverse phase counts a prefix of order \(\sqrt T(\log T)^2\). Differentiating its cubic remainder gains a factor of order the reciprocal rank in the spacing error as the rank grows. At a fixed positive time with an independently all-real divisor, the same even moments admit two reconstructions of a high zero. Direct polynomial evaluation with a circle count gives sufficient moment order and relative coefficient precision \(O_T(k+\log(1/\epsilon))\) for accuracy \(\epsilon\) in units of the \(k\)-th spacing. An imaginary-axis growth comparison identifies the absolute ranks in Polymath's eventual clock at \(T=1/2\).

For the time evolution we first express the gap velocity through a sum over all other zeros. We then tilt the complete source at the upper time. On an interior portion of the growing prefix the squared half-gap widens backward and has positive second derivative throughout a time interval of length \(\sqrt\kappa/(2L)\), where \(L\) is the upper saddle and \(\kappa\) its curvature. The leading curvature comes from the cubic asymmetry of the saddle. The remaining problems are to connect the height ranges at a specified time and to propagate counted gaps toward lower time.
\end{abstract}
\begingroup
\small
\tableofcontents
\endgroup
\clearpage
\section{The field, its zeros and the two limits}\label{sec:source}
Fix an upper physical time \(T\). We seek a formula for the \(k\)-th positive zero \(x_k(T)\) whose error can be compared with the spacing to its neighbor. A common error in two predicted positions can largely cancel in their difference. This makes the dependence of the error on \(k\) as important for gaps as its magnitude at a single root. An estimate for fixed \(k\) as \(T\to\infty\), an estimate uniform for \(k\le K(T)\), and an estimate as \(k\to\infty\) at one fixed \(T\) answer distinct questions.

In the normalization of Polymath \cite[\S1]{Polymath2019}, the physical field and its source are
\begin{align}
 \Phi(u)&=\sum_{n\ge1}(2\pi^2n^4 e^{9u}-3\pi n^2e^{5u})e^{-\pi n^2 e^{4u}},\nonumber\\
 H(t,z)&=\int_0^\infty e^{tu^2}\Phi(u)\cos(zu)\dd u,\qquad Z(t)=H(t,0).
 \label{eq:field}
\end{align}
Thus \(H(0,z)=\xi((1+\ii z)/2)/8\) and \(H_t=-H_{zz}\). Let \(U\ge0\) have probability law \(e^{tu^2}\Phi(u)\dd u/Z(t)\), and write \(a=a_t=\E_t U\).
The complex transform \(B(t,z)=Z(t)\E_t e^{\ii zU}\) packages the source cumulants. Its even part is \(H\), and \(H=\Real B\) on the real axis. The double-exponential tail permits differentiation and complex arguments in these formulas.

We count positive zeros from the origin. For a counted, simple, adjacent pair use
\[
 q_k=(k-\tfrac12)\pi,\quad x_-=x_k,\quad x_+=x_{k+1},\quad
 g_k=x_+-x_-,\quad c=(x_++x_-)/2,\quad D_k=g_k^2/4 .
\]
A dot denotes physical time \(t\). We suppress \(k\) when the pair is fixed. A descent of length \(s\) from \(T\) ends at \(T-s\).
The two limits require different information. At large \(T\) the source concentrates near one saddle. Its odd cumulants deform the zeros of a single cosine, and a common inverse preserves the dependence between neighbors (Section~\ref{sec:anchors}). At fixed \(T\) one can instead use the full even product and high-zero spacing to measure the precision needed to recover an absolute rank (Section~\ref{sec:comparison}).

For motion, knowing two positions leaves their relative velocity undetermined because every other zero enters it. Section~\ref{sec:exterior} completes that exterior from radial traces. Section~\ref{sec:anchored-curvature} obtains the signed second derivative from one upper source, including the influence of the unlocated zeros throughout the descent. Section~\ref{sec:current} derives a complementary velocity formula on intervals of increasing phase. The finite checks in Section~\ref{sec:finite} illustrate the error scales. Section~\ref{sec:frontier} states what would connect these results at a common height and extend them to lower time.
\section{Exact data at the large-time anchor}\label{sec:anchors}
The source and the half-integer lattice fix every quantity in the following expansions \cite{Founding07,Founding08}. At time \(t=T\), take the large positive saddle and its curvature
\begin{equation}
 4\pi e^{4L}=2TL+9,\qquad \kappa=8TL-2T+36,\qquad S=\kappa+2T,\qquad \ell=a\sqrt\kappa .
 \label{eq:saddle}
\end{equation}
In particular \(L=-W_{-1}(-8\pi e^{-18/T}/T)/4-9/(2T)\), \(L\sim(\log T)/4\), and \(\kappa\sim 8TL\). The normalized field at \(z=w/L\) tends to \(\cos w\) on compact sets \cite[positive-time saddle and Appendix A]{Saddle}.

\subsection{Evaluating the individual loops}
Put \(f(w)=H(T,w/L)/Z(T)\), \(d(w)=f(w)-\cos w\), \(q=q_k\), \(e=\sin q\), and \(r=1/2\). Require \(|d|<\sin r\) on the closed disks of radius \(r\) about the first \(k\) positive cosine zeros and on the uncovered parts of \([0,q+r]\). On each circle \(|\cos w|\ge\sin r\). Rouché and conjugation give one simple real zero there, and the real inequalities fix its index. These conditions hold for each fixed \(k\) eventually.
\begin{proposition}\label{prop:loop}
On the counterclockwise circle about \(q\), with \(\partial=\partial_w\),
\begin{equation}
 Lx_k-q=-\frac1{2\pi\ii}\oint\log(1+d/\cos w)\dd w
 =\sum_{n\ge1}\frac{e^n}{n!}\left[P_{n-1}(\partial)d(w)^n\right]_{w=q},
 \label{eq:loop}
\end{equation}
where \(P_0=1,\ P_1(X)=X,\ P_{n+1}=(X^2+n^2)P_{n-1}\). At \(\eta=\max|d|/\sin r<1\) on the circle the tail after term \(m\) is bounded by \(r\eta^{m+1}/[(m+1)(1-\eta)]\).
\end{proposition}
\begin{proof}
Integrate the difference of logarithmic derivatives times \(w\) and integrate by parts on the circle, where the ratio lies in the disk about one of radius \(\eta\). With \(h=w-q\) the denominator is \((-e\sin h)^n\). The first two residues of \(J(h)/\sin^n h\) are \(J(0)\) and \(J'(0)\). Twice integrating by parts in a residue and using
\((\partial_h^2+n^2)\csc^n h=n(n+1)\csc^{n+2}h\)
gives the recursion. Summing the absolute logarithm tail on the circle gives the stated bound. In particular the residue of \(e^{\ii v h}\csc^n h\) is \(i^{n-1}\prod_{j=1}^{n-1}(v+n-2j)/(n-1)!\).
\end{proof}
A first regrouping puts \(A(w)=\arcsin(e\,d(w))\) on its zero branch. The quotient \(f(q+h)/\cos(q+h)\) equals \(1-A(q+h)/h\) times a holomorphic unit. Indeed the unit is \(\cos((h+A)/2)\operatorname{sinc}((h-A)/2)/\operatorname{sinc} h\). Since \(\max|A|\le\arcsin(\max|d|)<r\), integrating the logarithm gives \(\sum_{m\ge0}[\partial^m A^{m+1}]_{q}/(m+1)!\) for the same displacement. Cauchy's estimate bounds its \(m\)-th term by \(r(\max|A|/r)^{m+1}/(m+1)\).

\subsection{The inverse series shared by the zeros}
Let \(\gamma_j=\partial_v^j\log\E_T e^{vU}|_{v=0}\) be the half-line cumulants. In a centered logarithm disk define
\begin{equation}
 K_0(y)=\log\E_T e^{\ii y(U-a)/a},\quad
 \Psi(y)=-\frac{K_0(y)-K_0(-y)}{2\ii}
      =\sum_{j\ge1}c_j y^{2j+1},\quad
 c_j=\frac{(-1)^{j+1}\gamma_{2j+1}}{(2j+1)!\,a^{2j+1}} .
 \label{eq:odd}
\end{equation}
Write \(y=az\) and \(\theta(y)=y-\Psi(y)\). The complex factorization on this disk is
\[
 \frac{H(T,z)}{Z(T)}=\exp\!\left(\frac{K_0(y)+K_0(-y)}2\right)
 \cos\theta(y),\qquad
 \cos\theta(y)=\prod_{k\ge1}\left(1-\frac{\theta(y)^2}{q_k^2}\right).
\]
The cosine product is classical \cite[4.22.2]{DLMF}. The amplitude has a logarithm on this disk. On the increasing real branch the individual zero solves \(\theta(y)=q_k\).

\begin{theorem}\label{thm:inverse}
The inverse \(Y\) through zero has the expansion
\begin{align}
 a x_k=Y(q_k),\qquad
 Y(q)&=\sum_{N\ge0}\frac{q^{2N+1}}{2N+1}
 [u^{2N}]\left(1-\sum_{j\ge1}c_j u^{2j}\right)^{-(2N+1)} .\label{eq:inverse}
\end{align}
In particular the coefficients of \(q^3,q^5,q^7\) are \(c_1,\ 3c_1^2+c_2,\ 12c_1^3+8c_1c_2+c_3\). For any fixed \(C>0\), the series is absolutely convergent for \(q_k\le C\ell\) at sufficiently large time depending on \(C\).
\end{theorem}
\begin{proof}
Apply Lagrange inversion to \(y-\Psi(y)=q\) \cite[\S1.10(vii)]{DLMF}. An effective disk test takes \(R>0\),
\[
 \delta=\frac{R^2}{2a^2}\E_T[(U-a)^2e^{R|U-a|/a}]<1,\quad
 b=-\log(1-\delta)/6,\quad Q=R/2-b>0 .
\]
The centered exponential Taylor remainder bounds \(|e^{K_0}-1|\le\delta\) on the disk of radius \(R\). Cauchy's estimate bounds \(\sum|c_j|(R/2)^{2j+1}\le b\). On \(|y|=R/2\), Rouché counts exactly one preimage for each \(|q|<Q\). It varies analytically and its real derivative is positive by continuation from zero. Replacing all \(c_j\) by their absolute values yields a majorant with displacement at most \(b\). Its tail after \(N=N_0\) is bounded by \(b(|q|/Q)^{2N_0+3}/(1-(|q|/Q)^2)\). For the larger asymptotic range use \eqref{eq:complex-saddle} below on \(|y|\le4C\ell\). The odd part has supremum \(O_C(\kappa^{-1/2})\) there. The same half-radius and majorant argument gives \(Q\ge2C\ell-O_C(\kappa^{-1/2})\).
\end{proof}
For \(h=-c_1>0\), all powers of the third cumulant alone sum explicitly to the branch
\[
 Y_3(q)=\frac{2}{\sqrt{3h}}\sinh\!\left(\frac13\operatorname{arsinh}
                          \frac{3\sqrt{3h}\,q}{2}\right),\qquad Y_3+hY_3^3=q .
\]
The remaining terms can be summed by the number of higher cumulants they contain. With \(I(y)=\Psi(y)-c_1y^3\), they give
\[
 Y=Y_3+\sum_{m\ge1}\frac{\partial_q^{m-1}}{m!}
                \{Y_3'(q)\,I(Y_3(q))^m\}.
\]
Insert a parameter \(\lambda\) multiplying \(I\) and invert at \(\lambda=0\). On every strict majorant disk above the one-root comparison persists for \(|\lambda|\le1+\epsilon\) at fixed \(q\), justifying evaluation at one.

The coefficients can themselves be expressed without quadrature. With \(b_n=\pi n^2\), \(J_b(v)=b^{-v}\Gamma(v,b)\), set
\[
 P(s)=\tfrac14\sum_{n\ge1}\!\left[2b_n^2J_{b_n}((9+s)/4)
               -3b_nJ_{b_n}((5+s)/4)\right],\qquad
 P_T(s)=\sum_{m\ge0}\frac{T^m P^{(2m)}(s)}{m!}.
\]
Substitute \(b_ne^{4u}\) in \eqref{eq:field} at \(T=0\), then expand the heat weight. The double-exponential tails justify absolute convergence on compact \(s,T\) sets, and \(\gamma_j=(\log P_T)^{(j)}(0)\) \cite[8.2.2]{DLMF}. An algebraic inverse germ in the continuous variable \(q\) is excluded at finite \(T\). Indeed put \(m(s)=\E_Te^{sU}\). At \(z=-\ii s\) the physical phase continues as \(-\ii s b(s)\), \(b(s)=\log(m(s)/m(-s))/(2s)\), because both transforms are positive there. Unbounded support gives \(b(s)\to\infty\); the density at zero gives \(m(-s)\gg_T1/s\), and the complete tail gives \(\log m(s)=O_T(s\log(2+s))\) by splitting at a sufficiently large multiple of \(\log(2+s)\). Hence \(b(s)=O_T(\log s)\). In a polynomial relation between argument and phase, divide by \(s^d\) at highest total degree \(d\). The leading homogeneous part is a nonzero polynomial in \(b(s)\), while every lower-degree term tends to zero. This contradiction concerns the continuous germ. A discrete-index or transcendental expression remains possible.

\subsection{The full set of factors}
At an all-real anchor the paired order-one product is \(H(T,z)/Z(T)=\prod_{i\ge1}(1-b_i z^2)\), \(b_i=x_i^{-2}\), counted with multiplicity \cite[\S2]{AGC}. Every individual factor can also be recovered as a limit from even moments \(m_j=\E_T U^j\), without a local inverse condition. Put
\[
 p_n=-n[v^n]\log\!\sum_{j\ge0}\frac{(-1)^jm_{2j}v^j}{(2j)!},\quad
 E_{0,n}=1,\quad E_{j,n}=\frac1j\sum_{h=1}^j(-1)^{h-1}p_{hn}E_{j-h,n}.
\]
Here \(p_n=\sum_i b_i^n\) and \(E_{j,n}\) sums the \(n\)-th powers of all products of \(j\) distinct indexed \(b_i\). Their sum at \(n=1\) is finite, so \(E_{j,n}^{1/n}\to b_1\cdots b_j\). Consequently \(x_k=\lim_n(E_{k-1,n}/E_{k,n})^{1/(2n)}\). High-order moment conditioning matters in using this identity.

For \(P_m(z)=\prod_{i\le m}(1-b_i z^2)\), Gaussian convolution gives the exact earlier field
\begin{equation}
 \frac{H(T-s,z)}{Z(T)}
   =\E_T[e^{-sU^2}\cos(zU)]
   =\lim_{m\to\infty}\sum_{j=0}^m\frac{s^j}{j!}\partial_z^{2j}P_m(z).
 \label{eq:whole}
\end{equation}
Indeed \(|P_m(w)|\le H(T,\ii|w|)/Z(T)\); its log growth is \(O_T(1+|w|\log(2+|w|))\). Dominated convergence against the translated Gaussian therefore holds locally uniformly. By contrast, the mass loss of the single frozen cosine with the exact initial curvature has the explicit size
\[
 \log\frac{Z(T-s)}{Z(T)e^{-s\E_T U^2}}
   =\int_0^s(s-v)\Var_{T-v}(U^2)\dd v .
\]
This follows by differentiating \(\log Z\) twice. Since \(\Var_T(U^2)\sim4L^2/\kappa\), the loss tends to one in logarithmic units at \(s=2\sqrt{T/L}\) \cite[\S4]{Founding07}. The local current below retains the variation responsible for actual gap motion.

\section{High ranks at a fixed all-real time}\label{sec:comparison}

Fix \(T>0\) at which the full divisor is known to be real. Set
\[
 F(z)=H(T,z)/Z(T),\qquad e_m=\frac{\E_T U^{2m}}{(2m)!},\qquad
 G(v)=\sum_{m\ge0}(-1)^m e_m v^m .
\]
Thus \(F(z)=G(z^2)\). We compare two uses of these same coefficients. The
first estimates an individual factor from symmetric functions. The second
counts polynomial zeros on circles where the full field has a margin
\cite{Founding07,EZHighRank}.

\subsection{The complete product and fixed-time spacing}

For \(u\ge0\), the source in \eqref{eq:field} is positive and is bounded
above by \(C e^{9u-\pi e^{4u}}\). Indeed, after factoring out the first
exponential, the sum of the upper bounds \(2\pi^2 n^4 e^{-\pi(n^2-1)}\)
is finite. For \(|z|\le r\), use \(|\cos(zu)|\le e^{ru}\).
Maximizing \(ru-\pi e^{4u}/2\) on the half line bounds it above by
\(O(r\log(2+r))\). The integral with the remaining factor
\(e^{Tu^2+9u-\pi e^{4u}/2}\) is finite. Hence the logarithm of the
maximum modulus is \(O_T(1+r\log(2+r))\), and Jensen's formula gives
\(O_T(r\log(2+r))\) zeros in the disk of radius \(r\) for \(r\ge2\).
In particular \(\sum_i x_i^{-2}<\infty\), counting positive zeros with
multiplicity. The entire function has order at most one, so Hadamard
factorization pairs the genus-one factors at \(\pm x_i\). Evenness removes
the linear exponential, and \(F(0)=1\) fixes the constant. Therefore
\begin{equation}\label{eq:all-real-product}
 F(z)=\prod_{i\ge1}(1-b_i z^2),\qquad
 G(v)=\prod_{i\ge1}(1-b_i v),\qquad b_i=x_i^{-2}.
\end{equation}
The products converge locally uniformly. In particular \(e_m\) is the
\(m\)-th elementary symmetric sum of the complete sequence \((b_i)\).

Ki, Kim and Lee's Theorems 1.3--1.4 give the simple real tail and its
asymptotic spacing at each fixed positive time \cite{KiKimLee2009}.
Their normalization is \(\Xi_\mu(w)=8H(4\mu,2w)\), so in physical
coordinates \(\Delta_i\log x_i\to4\pi\), where
\(\Delta_i=x_{i+1}-x_i\). The all-real assumption identifies this tail
with the positive product ranks up to a finite shift. Applying the
mean-value theorem to \(x\log x-x\) across each such gap gives increments
converging to \(4\pi\), since \(\Delta_i/x_i\to0\).
Telescoping and inverting gives
\begin{equation}\label{eq:high-rank-geometry}
 x_i\sim\frac{4\pi i}{\log i},\qquad
 \Delta_i\sim\frac{4\pi}{\log i},\qquad
 \sum_{i>j}b_i\ll_T\frac{\log^2(2j)}{j}.
\end{equation}
The tail-sum bound follows by comparison with
\(\int_j^\infty \log^2(2v)\,\dd v/v^2\). The onset of this spacing
regime at the chosen time is unevaluated.

\subsection{Recovering individual factors}

Let \(B_j=\prod_{i\le j} b_i\), \(B_0=1\), and let \(E_{j,h}\) be
the \(j\)-th elementary symmetric sum of the sequence \((b_i^h)_i\)
for integer \(h\ge1\). In particular \(E_{0,h}=1\). The logarithm
of \eqref{eq:all-real-product} near zero and Newton's identities give
\[
 p_n=-n[v^n]\log G(v)=\sum_i b_i^n,qquad
 E_{j,h}=\frac1j\sum_{n=1}^j(-1)^{n-1}p_{nh} E_{j-n,h} .
\]
The recurrence follows first for a finite head and then by absolute
convergence. Each \(E_{j,h}\) uses even moments through order \(2jh\).
Since \(B_j\) is the largest indexed product in \(E_{j,1}\),
\(B_j^h\le E_{j,h}\le B_j^{h-1}E_{j,1}\).
For \(h>1\) define
\(\underline B_{j,h}=(E_{j,h}/E_{j,1})^{1/(h-1)}\) and
\(\overline B_{j,h}=E_{j,h}^{1/h}\), with both equal to one at \(j=0\).
Then
\begin{equation}\label{eq:factor-interval}
 \sqrt{\frac{\underline B_{k-1,h}}{\overline B_{k,h}}}
 \le x_k\le
 \sqrt{\frac{\overline B_{k-1,h}}{\underline B_{k,h}}}.
\end{equation}
In particular \(x_k=\lim_{h\to\infty}(E_{k-1,h}/E_{k,h})^{1/(2h)}\)
for each fixed rank, including multiplicities. For finite input
intervals the alternating recurrence requires certified coefficient
errors before dividing by positive lower endpoints.

If the preceding factors are known, a separate estimate uses
\(s_h^{(k)}=p_h-\sum_{i<k}b_i^h\). The ratio
\(s_{h+1}^{(k)}/s_h^{(k)}\) is a weighted average of the factors
from index \(k\) onward. Subtract the average divided by \(b_k\)
from one, then use \(s_h^{(k)}\ge b_k^h\) and bound each
\((b_i/b_k)^{h-1}\) for \(i>k\) by the first such term. This gives
\[
 0\le 1-\frac{s_{h+1}^{(k)}}{b_k s_h^{(k)}}
 \le (b_{k+1}/b_k)^{h-1}\sum_{i>k} b_i/b_k .
\]
Strict separation gives geometric decay after the tail is bounded.
The subtraction also requires precision on the preceding factors
\cite[\S3]{Founding07}.

\paragraph{Exchanges across one rank.}
For \(b_{j+1}\le\theta_j\le b_j\), write
\[
 A_j=\sum_{i\le j}(\theta_j/b_i)^h,qquad
 C_j=\sum_{i>j}(b_i/\theta_j)^h .
\]
A term in \(E_{j,h}/B_j^h\) removes and inserts the same number of
indices across the cut at \(j\). All powers of \(\theta_j\)
therefore cancel. Writing \(e_r(\cdot)\) for an elementary symmetric
function here gives
\begin{equation}\label{eq:exchange-bound}
\begin{aligned}
 1\le \frac{E_{j,h}}{B_j^h}
 &=\sum_{r=0}^j e_r(((\theta_j/b_i)^h)_{i\le j})
                  e_r(((b_i/\theta_j)^h)_{i>j})\\
 &\le\sum_{r=0}^{\infty}\frac{(A_jC_j)^r}{(r!)^2}
 \le e^{A_jC_j}.
\end{aligned}
\end{equation}
The symmetric sums are bounded by the indicated powers divided by
\(r!\) on each side, by expanding the powers and counting the terms
with distinct indices. The tail sums converge by \(\sum_i b_i<\infty\).

If \(D_j\ge A_j C_j\) and \(0<L_j\le E_{j,h}\le U_j\), then
\begin{equation}\label{eq:exchange-interval}
 \left(\frac{L_{k-1}e^{-D_{k-1}}}{U_k}\right)^{1/(2h)}
 \le x_k\le
 \left(\frac{U_{k-1}e^{D_k}}{L_k}\right)^{1/(2h)} .
\end{equation}
Use \(D_0=0\) and \(L_0=U_0=1\).
Indeed, \(B_j^h\) lies between \(L_j e^{-D_j}\) and \(U_j\), and
\(x_k^{2h}=B_{k-1}^h/B_k^h\). The logarithm of the ratio of its two
endpoints is
\[
 \frac{D_{k-1}+D_k+\log(U_k/L_k)+\log(U_{k-1}/L_{k-1})}{2h}.
\]
A finite counted head through \(J\ge j+1\) supplies \(D_j\) if a
threshold can be chosen between the upper bound for \(b_{j+1}\) and
the lower bound for \(b_j\). Bound the known weights by interval
endpoints. For the rest, if \(b_{J+1}\le\beta\le\theta_j\) and
\(\sum_{i>J}b_i\le S\), the remaining contribution to \(C_j\)
is at most \((\beta/\theta_j)^{h-1}S/\theta_j\). The upper
endpoint for \(b_J\) can supply \(\beta\), and
\(\sum_i b_i=e_1\) supplies \(S\) after subtracting a lower
bound on the head sum. If both nearest threshold ratios are at
most \(q<1\), then \(A_j\le jq^h\) and
\(C_j\le q^{h-1}\sum_{i>j}b_i/\theta_j\). This gives geometric
decay for fixed \(j\). At a repeated factor \eqref{eq:factor-interval}
still applies.

At high ranks the spacing gives a uniform allowance for each fixed
\(T\). Take \(\theta_j=(b_j b_{j+1})^{1/2}\).
For \(\lfloor j/2\rfloor\le i\le 2j\), the adjacent differences
\(\log b_i-\log b_{i+1}=2\log(1+\Delta_i/x_i)\) are at least
\(c_T/j\) by \eqref{eq:high-rank-geometry}. The nearest weights on
each side of the cut thus cost half a step, and each subsequent
weight at least one additional step. On this range either sum is at
most \(e^{-c_T h/(2j)}/(1-e^{-c_T h/j})\).
On the far head, each base \(\theta_j/b_i\) is at most a fixed
\(q_0<1\) since \(x_{\lfloor j/2\rfloor}/x_j\to1/2\).
On the far tail each base \(b_i/\theta_j\) is also at most
\(q_0<1\) after increasing \(q_0\), since \(x_{2j}/x_j\to2\).
The head costs at most \(j q_0^h\), and the tail costs at most
\(q_0^{h-1}\theta_j^{-1}\sum_{i>2j}b_i=O_T(jq_0^{h-1})\).
For \(h\ge j\) these far bounds and both near sums are
\(O_T(e^{-c'_T h/j})\) with some \(c'_T>0\). For example,
one can absorb the factor \(j\) into half the exponential decay
in \(h\). We can therefore take \(D_j=C_Te^{-c_T'' h/j}\)
in \eqref{eq:exchange-interval} for sufficiently large \(j\)
and all \(h\ge j\).

\paragraph{Uncertainty in the symmetric coefficients.}
For \(\omega=e^{2\pi\ii/h}\), multiplying the rotated products gives
\[
 \prod_{l=0}^{h-1}G(\omega^l v)=\prod_i(1-b_i^h v^h),
 \qquad
 (-1)^j E_{j,h}=[v^{jh}]\prod_{l=0}^{h-1}G(\omega^l v).
\]
The product is locally absolutely convergent, and for \(j\le k\)
only the coefficients through \(kh\) are used. Equip these
truncated polynomials with
\(\|P\|_R=\sum_m |[v^m]P|R^m\), for \(R>0\).
The norm is submultiplicative and rotations are isometries.
Let \(\widetilde G\) be the center polynomial and \(d_R\) the
weighted sum of its coefficient radii in this norm. Put
\(U_R=\|\widetilde G\|_R+d_R\). Replacing the rotated factors
one at a time gives a telescoping sum. Taking its norm and then
extracting degree \(jh\) gives
\begin{equation}\label{eq:exchange-coefficient-error}
 |\widetilde E_{j,h}-E_{j,h}|
 \le h U_R^{h-1}d_R/R^{jh}.
\end{equation}
Enclose any arithmetic error separately.

Write \(M_R=G(-R)=F(\ii\sqrt R)\). The exact truncated norm is
at most \(M_R\). If \(d_R\le\eta M_R\), then
\(U_R\le(1+2\eta)M_R\). Since \(E_{j,h}\ge B_j^h\),
the relative error in \eqref{eq:exchange-coefficient-error} is at most
\[
 h\eta(1+2\eta)^{h-1}
 \left(\frac{M_R}{R^j B_j}\right)^h.
\]
Take \(R=(4\pi k/\log k)^2\). For \(j=k-1,k\),
\begin{equation}\label{eq:exchange-coefficient-growth}
 \frac{M_R}{R^j B_j}
 =\prod_{i\le j}(1+x_i^2/R)\prod_{i>j}(1+R/x_i^2)
 \le e^{C_T k}.
\end{equation}
Every head factor is bounded by a constant since \(x_j^2\asymp R\).
The tail logarithm is at most \(R\sum_{i>j}b_i=O_T(k)\).
Thus coefficient radii \(2^{-p} e_m\) through degree \(kh\)
are sufficient for relative output errors \(O(\tau)\) if
\(p\ge C_T'kh+\log_2(h/\tau)+C_T''\), with constants large enough
to absorb the factor \((1+2\eta)^{h-1}\) and the passage from
radii to endpoints.

\begin{theorem}[Sufficient product input budget]\label{thm:high-rank-exchange}
For each fixed positive all-real \(T\), at sufficiently large rank
\(k\), \(h\ge k\) and relative coefficient enclosure radii at most
\(\tau\le1/4\) for \(E_{k-1,h},E_{k,h}\), the width of
\eqref{eq:exchange-interval} in units of \(\Delta_k\) is
\[
 O_T\!\left(\frac{k}{h}(e^{-c_T h/k}+\tau)\right).
\]
In particular, \(h\asymp_T k(1+\log(1/\varepsilon))\) with a
sufficiently large constant and \(\tau\asymp\varepsilon\)
give width \(O_T(\varepsilon\Delta_k)\). The sufficient even
moment order is \(2kh\), and the sufficient input precision per
normalized coefficient is \(O_T(kh+\log(h/\tau))\) bits.
\end{theorem}
\begin{proof}
Substitute the high-rank bound for \(D_j\) into the logarithmic
width of \eqref{eq:exchange-interval}. Each log coefficient
ratio is \(O(\tau)\). The logarithmic width is \(O_T(1/h)\),
so exponentiating and using \(x_k/\Delta_k\asymp k\) proves
the width estimate. The precision bound follows from
\eqref{eq:exchange-coefficient-error}--\eqref{eq:exchange-coefficient-growth}.
\end{proof}
The sufficient asymptotic constants come from the spacing theorem.
At a specified rank one can use the certified head and tail
allowance above. The symmetric coefficients themselves are evaluated
directly from the moments.


\subsection{Direct evaluation from the moments}

Use the polynomial \(P_N(z)=\sum_{m=0}^N(-1)^m e_m z^{2m}\).
For real \(x\), Taylor's theorem through order \(2N+1\) inside the
positive cosine integral gives
\begin{equation}\label{eq:moment-real-error}
 |F(x)-P_N(x)|\le e_{N+1}|x|^{2N+2}.
\end{equation}
If \(\widetilde P_N\) uses coefficient centers with radii \(d_m\), add
\(\sum_{m=0}^N d_m X^{2m}\) on \(|x|\le X\) and use an upper
endpoint for \(e_{N+1}\). A numerical evaluation must also enclose its
arithmetic error.

The defining source gives refinable coefficient intervals at any fixed
\(T\). For example, let \(J_m=\int_0^\infty u^{2m}e^{Tu^2}\Phi(u)\dd u\).
To truncate at \(U_0\ge1\), set
\[
 p_m(u)=Tu^2+9u+2m\log u-\pi e^{4u},\qquad
 A=-p_m'(U_0)>0,qquad 16\pi e^{4U_0}\ge2T .
\]
These conditions hold for all sufficiently large \(U_0\).
Concavity on \([U_0,\infty)\) gives
\(p_m(u)\le p_m(U_0)-A(u-U_0)\). Hence the omitted integral is
at most \(2\pi^2 C_\Phi e^{p_m(U_0)}/A\), where
\(C_\Phi=\sum_{n\ge1}n^4 e^{-\pi(n^2-1)}\).
On \([0,U_0]\), truncating the square modes at \(n=M\) costs at most
\[
 2\pi^2 U_0^{2m+1} e^{TU_0^2+9U_0}
       \sum_{n>M}n^4 e^{-\pi(n^2-1)} .
\]
The displayed sums admit geometric upper bounds, for instance by
\(n^2-1\ge n-1\) and the power series for \(\sum n^4 q^{n-1}\) at
\(q=e^{-\pi}\). On the remaining compact integral, upper and lower
interval sums converge as the mesh tends to zero since the finite
integrand is continuously differentiable. Divide the resulting
enclosures as \(e_m=J_m/((2m)!J_0)\) using a positive lower endpoint
for \(J_0\). This gives a finite error allowance for generation of the
inputs. The order and precision estimates below concern the moments
supplied to the polynomial.

\begin{lemma}[Field bounds on circles near a high zero]\label{lem:high-rank-circle}
For every \(r\ge0\), \(\min_{|z|=r}|F(z)|=|F(r)|\).
There is a fixed \(d>0\) such that, for all sufficiently large
\(k\) and each fixed \(A>0\),
\begin{equation}\label{eq:direct-growth}
\begin{aligned}
 |F(r)|&\ge e^{-C_T k}\frac{|r-x_k|}{k\Delta_k}
             && (|r-x_k|\le d\Delta_k),\\
 F(2\ii X)&\le e^{C_{T,A}k}
             && (0\le X\le A k/\log k).
\end{aligned}
\end{equation}
\end{lemma}
\begin{proof}
On \(|z|=r\), the modulus of each factor in \eqref{eq:all-real-product}
is at least \(|1-b_i r^2|\), with equality for real \(z\). This
passes to the full product. The factor at \(i=k\) has magnitude
\(|r-x_k|(r+x_k)/x_k^2\gg_T |r-x_k|/(k\Delta_k)\).
For \(\lfloor k/2\rfloor\le i\le2k\), the gaps are uniformly at
least a fixed multiple of \(\Delta_k\) by \eqref{eq:high-rank-geometry}.
Take \(d\) below half this multiple. Then for \(i\ne k\),
\(|r-x_i|\gg_T |i-k|\Delta_k\) and
\(|1-r^2/x_i^2|\ge c_T |i-k|/k\), using \(x_i\ll_T k\Delta_k\)
and \(r+x_i\ge x_i\). The product of these bounds is
\(\ge e^{-C_T k}\) by \(m!\ge(m/e)^m\) on the two sides of \(k\).
Below \(\lfloor k/2\rfloor\) each factor has magnitude at least one
eventually, since \(x_{\lfloor k/2\rfloor}/x_k\to1/2\).
Above \(2k\), the quantities \(r^2/x_i^2\) are bounded away
from one since \(x_{2k}/x_k\to2\). Their sum is \(O_T(k)\) by
\eqref{eq:high-rank-geometry}, so the logarithm of the tail product
is bounded below by \(-C_T k\).

For the last assertion, split
\(\log F(2\ii X)=\sum_i\log(1+4X^2/x_i^2)\) at \(2k\).
The head satisfies \(X/x_i\le C_{T,A}k/i\) there by
\eqref{eq:high-rank-geometry}, with the comparison constant enlarged
to cover the finite prefix. For \(i\le2k\), \(1+4X^2/x_i^2\le
C_{T,A}'(2k/i)^2\). The head logarithm is \(O_{T,A}(k)\) by the
factorial bound. The tail logarithm is at most
\(4X^2\sum_{i>2k} b_i=O_{T,A}(k)\).
\end{proof}

\paragraph{Finite polynomial count.}
On \(|z|\le X\) a bound on the exact omitted tail is
\(4^{-(N+1)}F(2\ii X)\), obtained by extracting \(4^{-(N+1)}\)
from the positive coefficient sum at \(2\ii X\). There is also an
a posteriori test using a single extra coefficient. The \(e_m\)
are log-concave. To see this, for a finite product
\(f(v)=\prod_i(1+b_i v)\), let \(p=f^{(m-1)}\). Rolle's theorem
and the real roots of \(p\) give \((p')^2-pp''\ge0\), by
differentiating \(p'/p\) away from the roots and then using continuity.
At zero this gives \((m!e_m)^2\ge (m-1)!(m+1)!e_{m-1}e_{m+1}\)
for the finite coefficients. Pass to the locally uniform product limit.
For \(N\ge1\), the decreasing ratios yield
\begin{equation}\label{eq:moment-circle-error}
 \sup_{|z|\le X}|F(z)-\widetilde P_N(z)|
 \le \sum_{m=0}^N d_m X^{2m}
       +\frac{\overline e_{N+1}X^{2N+2}}{1-\bar\eta},
 \qquad
 \bar\eta=\frac{X^2\overline e_{N+1}}{\underline e_N}<1 ,
\end{equation}
where the lower coefficient endpoint must be positive. To justify a
strict comparison at a circle of radius \(r\le X\), bound
\(\min_{|z|=r}|\widetilde P_N(z)|\) below by a value larger than
\eqref{eq:moment-circle-error}. This can be certified by covering the
circle with sufficiently small arcs and bounding variation using the
polynomial derivative. Rouch\'e transfers its root count to \(F\).
The polynomial count can be determined by finite winding arithmetic
on the same circle. Its value is twice the number of positive zeros
inside, with multiplicity, since the full divisor is real and even.
Counts \(2(k-1)\) and \(2k\) on two radii place the simple \(x_k\)
strictly between them.

\begin{theorem}[Sufficient direct input budget]\label{thm:high-rank-direct}
At each fixed positive all-real \(T\), for sufficiently large \(k\) and
\(0<\varepsilon\le1/4\), the finite polynomial test can bracket
\(x_k\) to width \(O_T(\varepsilon\Delta_k)\) using
\(O_T(k+\log(1/\varepsilon))\) as both the maximum even moment order
and a sufficient relative precision in bits per normalized coefficient.
\end{theorem}
\begin{proof}
Choose \(X=8\pi k/\log k\) so that \(x_k\) is interior to the search
interval \([X/4,X]\) for all sufficiently large \(k\). A grid
of mesh \(c\varepsilon/\log k\) with small fixed \(c>0\) has points
on both sides of \(x_k\) between one and three mesh widths away
and within \(d\Delta_k\). At each of those radii the minimum
modulus of \(F\) is at least \(e^{-C_T k}\varepsilon/(C_T k)\)
by Lemma~\ref{lem:high-rank-circle}. The two exact counts there
are \(2(k-1)\) and \(2k\).

The true ratio in \eqref{eq:moment-circle-error} satisfies
\(X^2 e_{N+1}/e_N\le(e_N X^{2N})^{1/N}
\le 4^{-1}\exp(C_T k/N)\), because the earlier ratios have
product \(e_N X^{2N}\) and \(e_0=1\). For \(N\ge C'_T k\)
with sufficiently large \(C'_T\), small relative coefficient radii
keep \(\bar\eta\) bounded strictly below one. The resulting
geometric tail bound is \(O(4^{-(N+1)}e^{C_T k})\).
Uniform relative radii \(\delta\) in the lower coefficients
contribute at most \(\delta F(\ii X)\le\delta e^{C_T k}\).
Taking \(N\) and \(\log(1/\delta)\) each a sufficiently large
multiple of \(k+\log(1/\varepsilon)\) makes the sum smaller than a
quarter of the circle margin. The polynomial minimum therefore
exceeds the sum at both chosen radii, and sufficiently refined
interval arithmetic certifies these strict tests. The input order
is \(2N+2\).
\end{proof}

The finite test also applies at a specified lower rank without a
spacing theorem. Enumerate rational radii, moment orders and
precisions, interleaving the boundary checks. For a fixed radius
avoiding the divisor, compact convergence gives a positive margin
eventually. The ratio bound eventually succeeds because the series
is entire with positive coefficients and decreasing ratios.
Thus every specified simple root can be bracketed to arbitrarily
small width. Two such adjacent brackets give a positive lower
bound on their separation and can be refined to a prescribed
fraction of that bound. The theorem gives an asymptotic input budget
and leaves search, polynomial arithmetic and moment generation costs
to be quantified.

\subsection{Matching the absolute clock rank at \texorpdfstring{\(t=1/2\)}{t=1/2}}

For \(0<t\le1/2\), Polymath's Theorem~1.3 bounds the error of
the normalized approximation to \(H(t,x+\ii y)\) for \(x\ge200\),
\(0\le y\le1\) \cite{Polymath2019}. Its high-zero theorem uses
\[
 G_t(x)=\frac{x}{4\pi}\left(\log\frac{x}{4\pi}-1\right)
           +\frac{11}{8}+\frac{t}{16}\log\frac{x}{4\pi},
 \qquad G_t(u_j)=j .
\]
Theorem~1.5(i)--(ii) gives one simple real zero \(r_j\) in each
disk of radius \(c/\log u_j\) for \(j\ge\exp(C/t)\). Every zero
with real part \(x\ge\exp(C/t)\) is real and is within
\(O(x^{-ct})\) of a clock center.
Here \(C\) is a sufficiently large absolute constant and
\(c>0\) a sufficiently small one.
The \href{https://arxiv.org/html/1904.12438v2\#S9.E104}{proof after (104)} gives
\begin{equation}\label{eq:clock-error}
 |r_j-u_j|\ll\frac{u_j^{-ct}}{\log u_j},\qquad
 s_j=u_{j+1}-u_j\sim\frac{4\pi}{\log(u_j/(4\pi))}.
\end{equation}
The spacing follows by differentiation of \(G_t\). Theorem~1.5(iii)
gives \(N_t(X)=G_t(X)+O(1)\), where \(N_t\) counts zeros from the
imaginary axis by real part, with multiplicity. At an all-real
time this count is the positive rank from the origin
\cite[Theorem~1.5 and \S9]{Polymath2019}.

\begin{proposition}[Absolute clock rank]\label{prop:clock-rank}
At \(t=1/2\), \(r_j=x_j\) for all sufficiently large \(j\).
\end{proposition}
\begin{proof}
The divisor at \(t=1/2\) is all real by de Bruijn's upper bound
\cite[\S1]{Polymath2019}. The windows in \eqref{eq:clock-error}
are disjoint and ordered for large \(j\). Consequently for
\(x\ge b\), with \(b\) fixed sufficiently large and in a gap
between the windows, one has
\[
 N_t(x)=\lfloor G_t(x)\rfloor+d_0+\rho(x)
\]
for an integer \(d_0\). The error satisfies \(|\rho|\le1\), and
apart from interval endpoints its support lies between each clock
center and the corresponding root. Each dyadic interval \([R,2R]\)
contains \(O(R\log R)\) clock centers. By \eqref{eq:clock-error}
the integral of \(|\rho(x)|/x\) on that interval is
\(O(R^{-ct})\). Thus \(\rho\) is integrable against \(\dd x/x\)
on \([b,\infty)\).

On the imaginary axis use the complete product at time \(t\).
For \(Y>0\), integrate
\(K_Y(x)=2Y^2/[x(x^2+Y^2)]\) above each positive root.
The integral for root \(x_i\) is \(\log(1+Y^2/x_i^2)\).
Positivity permits summation under the integral, giving
\[
 \log\frac{H(t,\ii Y)}{Z(t)}
   =\int_0^\infty K_Y(x)N_t(x)\dd x .
\]
The finite head contributes \(O(1)\) as \(Y\to\infty\), since
\(K_Y(x)\le2/x\) and there is no zero near the origin.
The error \(\rho\) also costs \(O(1)\).
The function \(p(s)=\lfloor s\rfloor-s+1/2\) has a bounded
primitive by its zero mean on each period. Since \(G_t'\)
is eventually positive and \(G_t''(x)=1/(4\pi x)-t/(16x^2)\),
we may choose \(b\) so that \(G_t'\) is increasing there.
With \(s=G_t(x)\), the weight for \(p(s)\)
is \(K_Y(x)/G_t'(x)\), decreasing to zero and bounded at
\(b\) uniformly in \(Y\). Integration by parts shows that
this term is \(O(1)\). It remains to integrate
\(G_t(x)+d_0-1/2\) over \([b,\infty)\).

The term in \(x(\log(x/(4\pi))-1)/(4\pi)\) can be extended
down to zero with an \(O(1)\) error. Substitute \(x=Yu\).
The integral of \((1+u^2)^{-1}\) is \(\pi/2\), and its
integral against \(\log u\) vanishes by \(u\mapsto1/u\).
The resulting contribution is
\(Y(\log(Y/(4\pi))-1)/4\). The constant \(7/8+d_0\)
is multiplied by \(\log(1+Y^2/b^2)=2\log Y+O(1)\).
For the remaining logarithmic term, replacing \(K_Y(x)\)
by \(2\mathbf 1_{x\le Y}/x\) gives \(\log^2(Y/(4\pi))+O(1)\)
as its integral against \(\log(x/(4\pi))\). The difference
uses the weight
\[
 W(u)=\frac2u\left(\frac{1}{1+u^2}-\mathbf 1_{u\le1}\right)
\]
in the scaled variable. Its integral is zero by inversion,
its absolute logarithmic moment is finite, and its integral
in absolute value on \([0,b/Y]\) is \(O(Y^{-2})\).
Thus the replacement costs \(O(1)\) even after multiplication
by \(\log(Y/(4\pi))+\log u\). Altogether,
\begin{equation}\label{eq:clock-count-growth}
 \log H(t,\ii Y)
 =\frac{Y}{4}\left(\log\frac{Y}{4\pi}-1\right)
   +\frac{t}{16}\log^2\frac{Y}{4\pi}
   +\left(\frac74+2d_0\right)\log Y+O(1).
\end{equation}

To evaluate the same logarithm from the source, put
\(F_0(Y)=H(0,\ii Y)\).
By evenness and the zero-time normalization,
\(F_0(Y)=\xi((1+Y)/2)/8\) for \(Y>0\). At
\(s=(1+Y)/2\), Stirling's formula and \(\zeta(s)=1+o(1)\)
give \((s/2)\log(s/(2\pi))-s/2+(3/2)\log s+O(1)\)
for \(\log F_0(Y)\). Substituting \(s\) gives
\[
 \log F_0(Y)=P(Y)+O(1),\qquad
 P(Y)=\frac{Y}{4}\left(\log\frac{Y}{4\pi}-1\right)
                       +\frac74\log Y .
\]
In particular the shift \(Y+1\) contributes \(\tfrac14\log Y\).
For a centered Gaussian \(V\) of variance \(2t\), the positive
source integral gives \(H(t,\ii Y)=\E F_0(Y+V)\).
Indeed the Gaussian average of \(\cosh((Y+V)u)\) is
\(e^{tu^2}\cosh(Yu)\), with interchange justified by positivity.

Write \(A_Y=P'(Y)=\tfrac14\log(Y/(4\pi))+7/(4Y)\).
On \(|v|\le Y^{1/4}\), Taylor's theorem gives
\(P(Y+v)=P(Y)+A_Yv+O(Y^{-1/2})\). Multiplying the
Gaussian density by \(e^{A_Yv}\) gives a shifted density of
mean \(2tA_Y=O(\log Y)\) times the factor \(e^{tA_Y^2}\).
The shifted Gaussian probability of the central interval tends
to one. On \(Y^{1/4}<|v|\le Y/2\), the bound
\(F_0(Y+v)\ll e^{P(Y)+C|v|\log Y}\) follows from the
estimate for \(P'\). The quadratic Gaussian exponent
dominates the linear term and gives a contribution of
\(O(e^{P(Y)-c_tY^{1/2}})\). On \(|v|>Y/2\),
evenness and the Stirling bound give
\(\log F_0(Y+v)=O((1+|v|)\log(2+|v|))\) since
\(|Y+v|\le3|v|\). This range contributes
\(O(e^{-c_tY^2})\). It follows that
\[
 \log H(t,\ii Y)=P(Y)+tA_Y^2+O(1)
               =P(Y)+\frac{t}{16}\log^2\frac{Y}{4\pi}+O(1).
\]
Comparison with \eqref{eq:clock-count-growth} forces \(d_0=0\).
The point with \(G_t(X)=j+1/2\) is eventually outside the
windows and has exactly \(j\) positive zeros through it.
This gives the asserted rank.
\end{proof}

The full published displacement at these matched ranks gives
position and gap errors \(O(u_k^{-ct})\) in units of \(s_k\).
For the gap, subtract the two displacements in
\eqref{eq:clock-error} and use \(u_{k+1}/u_k\to1\) and the
displayed spacing. This calibration leaves the height onset
unevaluated. A finite alternative at a cutoff in a proved clock
gap is an exact prefix count, for example by the finite
polynomial count above. If the adjacent clock labels there are
\(j,j+1\), subtract \(j\) from that count to obtain the offset.

For refinement with moments at \(t=1/2\), the matched clock
gives a short real interval containing just \(x_k\), for
example \([u_k-d's_k,u_k+d's_k]\) for small fixed \(d'>0\).
By \eqref{eq:clock-error} the whole interval eventually lies
within the neighborhood in Lemma~\ref{lem:high-rank-circle}.
Its endpoints have opposite signs and magnitudes
\(\gg e^{-Ck}/k\). With coefficient and arithmetic errors
small compared with \(e^{-Ck}\varepsilon/k\), the real error
bound \eqref{eq:moment-real-error} can be used to bisect this
interval. If a test interval contains zero, the modulus of the
true field is at most twice the certified radius, and the
lower bound in Lemma~\ref{lem:high-rank-circle} puts the test
point within \(O(\varepsilon s_k)\) of the root. Otherwise the
sign selects the next half-interval. The same moment order and
sufficient precision as in Theorem~\ref{thm:high-rank-direct}
meet this error allowance by \eqref{eq:direct-growth}.

With the same coarse geometry, Theorem~\ref{thm:high-rank-direct}
uses a shorter moment prefix at large rank than the sufficient
budget in Theorem~\ref{thm:high-rank-exchange}. If the full
clock displacement is available, it already supplies the
displayed localization accuracy at \(t=1/2\) without acquiring
moments. For an additional prescribed refinement the theorems
give sufficient orders and precision with either construction.
Subtracting the neighboring root intervals gives gap uncertainty
\(O_T(\varepsilon\Delta_k)\) since successive spacings are
asymptotically equal.
They leave the onset and arithmetic work unevaluated, as well as
the height overlap with the large-time inverse in
Section~\ref{sec:anchors} at a common \(T\).

\subsection{The published time-zero benchmark}

The identity \(H(0,z)=\xi((1+\ii z)/2)/8\) gives \(x=2\gamma\)
for a critical-line zero of ordinate \(\gamma\). Platt and
Trudgian certify that the first \(12\,363\,153\,437\,138\) zeros
of positive ordinate through \(\gamma=3\,000\,175\,332\,800\)
lie on the line, corresponding to \(x=6\,000\,350\,665\,600\)
\cite[Theorem~1 and \S2]{PlattTrudgian2021}. Their sign-change
count with Turing's method certifies completeness at time zero.
A positive-time count at a specified vertical boundary also
requires accounting for boundary crossings or counting at that time.

Polymath proved \(\Lambda\le0.22\)
\cite[Theorem~1.1]{Polymath2019}. The second row of its
Table~1 uses an ordinate verification above \(2.51\cdot10^{12}\),
so the larger Platt--Trudgian verification gives \(\Lambda\le0.2\)
by \href{https://arxiv.org/pdf/2004.09765v1\#page=5}{Corollary~2}
of \cite{PlattTrudgian2021}. The criterion giving
\(\Lambda\le t_0+y_0^2/2\) combines a time-zero verification
through ordinate \(X/2\), final-time nonvanishing for
\(x\ge X+\sqrt{1-y_0^2}\) and \(y_0\le y\le\sqrt{1-2t_0}\),
and an intermediate-time barrier on
\(X\le x\le X+\sqrt{1-y_0^2}\) above
\(\sqrt{y_0^2+2(t_0-t)}\) up to \(\sqrt{1-2t}\)
\cite[Theorem~1.2]{Polymath2019}. An improvement by this
criterion requires these conditions at parameters satisfying
\(t_0+y_0^2/2<0.2\). The factor and circle estimates concern
one independently established all-real time.

\section{From zero locations to the exterior interaction}\label{sec:exterior}

At a time \(T\) with a known all-real divisor, let \(x_-=x_k\) and
\(x_+=x_{k+1}\) be consecutive simple positive zeros. Write
\(g_k=x_+-x_-\) and \(D_k=g_k^2/4\). For integers \(p,q\ge0\) with
\(m=p+q\ge2\), the parent hierarchy uses
\[
 S_{k,p,q}=\sum_{\rho\ {\rm exterior}}
       (\rho-x_+)^{-p}(\rho-x_-)^{-q}.
\]
The sum counts multiplicities and omits just the two marked occurrences.
Thus the denominators are the displacements from the midpoint shifted by
the two signs of \(\sqrt{D_k}\). In the parent's notation
\(S_{k,p,q}=P^{(k)}_{p,q}\). On a separated real path the time derivative
of each entry uses the gap and finitely many entries of the same array,
including higher orders \cite[\S3, integer split recurrence]{AGC}.
For example, suppressing \(k\) on the right and left, an axis entry satisfies
\[
 S_{2,0}'=2\left(S_{2,0}^2-5S_{4,0}-2S_{3,1}+\frac{2}{g_k}S_{3,0}\right).
\]
The array is closed under differentiation when all orders and the gap are retained.
At a simple zero \(x_i\), differentiating \(H(t,x_i(t))=0\) gives
\(\dot x_i=H_{zz}/H_z\). Deflate its linear factor and differentiate
the paired product \eqref{eq:all-real-product}. This gives twice the
sum of \((x_i-\rho)^{-1}\) over the remaining zeros, paired by reflection.
In the difference of the two velocities the marked endpoints contribute
\(4/g_k\). Each common exterior occurrence contributes
\(-2g_k/((\rho-x_+)(\rho-x_-))\), and this difference is absolutely
summable. Consequently
\begin{equation}\label{eq:exterior-gap}
 M_k=S_{k,1,1},\qquad
 g_k'=\frac{4}{g_k}-2g_k M_k,\qquad D_k'=2-g_k^2M_k .
\end{equation}

Write \(Q_\rho=(\rho-c)^2-D=(\rho-x_-)(\rho-x_+)\) and
\(S_j=S_{k,j,j}=\sum_{\rho\ {\rm exterior}}Q_\rho^{-j}\) for the
diagonal traces. Thus \(M=S_1\) when the pair index is suppressed.
Differentiating the first mixed entry by the parent hierarchy
\cite[\S3, first diagonal row]{AGC} gives
\begin{equation}
 \dot D=2-4DM,\qquad
 \dot M=2M^2-\mathcal F,\qquad \mathcal F=12S_2+16D S_3 .
 \label{eq:parent}
\end{equation}
In particular the first derivative of the gap requires \(M\) and its
second derivative already involves higher exterior orders. We now give
a way to enclose these quantities from counted positions and source
moments at one time.

\paragraph{Completing the exterior from a counted head.}
Suppose the first \(K\ge k+1\) positive zeros are counted, and
\(x_+\le A<R\le x_i\) for every \(i>K\). Sum the contributions of the
first \(K\) reflected pairs directly, retaining \(-x_-\) and \(-x_+\).
For interval enclosures, each difference used with a positive exponent
must have a positive lower bound on its absolute value after allowing for
errors in both zeros. For the remaining pairs put
\[
 W_j=\sum_{i>K}x_i^{-2j},\qquad
 h_d=[v^d](1-x_+v)^{-p}(1-x_-v)^{-q}.
\]
The even product gives \(W_j\) by subtracting the counted head from the
full radial trace \(p_j\). In particular
\(W_1=\E_T U^2/2-\sum_{i\le K}x_i^{-2}\)
\cite{Founding07,EZMixedAnchor}. The complete tail of \(S_{k,p,q}\) equals
\begin{equation}\label{eq:exterior-tail}
 2\sum_{\substack{d\ge0\\m+d\ {\rm even}}} h_d W_{(m+d)/2}.
\end{equation}
Indeed reflection cancels terms of odd total degree. Since
\(0\le h_d\le \binom{m+d-1}{d} A^d\), for \(A/R<\sigma<1\)
the terms with \(d\ge L\) contribute at most
\[
  \frac{2R^{2-m} W_1}{(1-\sigma)^m}
       \left(\frac{A}{R\sigma}\right)^L .
\]
This bounds the truncation for any fixed row, including the axes.
An enclosure must also propagate the widths of the radial traces and
the head through their subtraction. For \(M_k\) the first tail term is
\(2W_1\).

\paragraph{A signed large-time readout.}
Set \(y_i=a x_i\), \(q_i=(i-\tfrac12)\pi\), and
\(\ell=a\sqrt{\kappa}\). Take \(K\asymp\ell\) inside a fixed multiple
of the large-time inverse window \cite{Founding08}. At fixed \(p,q\)
and \(0<\theta<1\), for \(k+1\le\theta K\) the counted separations and
the absolute tail bound give \cite{EZMixedAnchor}
\[
 a^{-m} S_{k,p,q}
 =\pi^{-m}\sum_{r\in\mathbb Z\setminus\{0,1\}}(r-1)^{-p}r^{-q}
    + O\!\left((\ell\sqrt{\kappa})^{-1}+K^{1-m}\right).
\]
For \(p=q=1\) the leading lattice interaction cancels the constant in
\eqref{eq:exterior-gap}. The more precise radial comparison of
\cite{EZMixedAnchor} gives, as \(T\to\infty\), uniformly for
\(k+1\le\sqrt{\ell}\),
\begin{equation}\label{eq:exterior-signed}
 \ell^2(g_k^2M_k-2)
   =\pi^2+O\!\left(\kappa^{-1/2}+\frac{(k+1)^2}{\ell}\right).
\end{equation}
The constants depend on the fixed window and the bounds on \(K/\ell\),
and in the general row estimate also on \(p,q,\theta\).
On the comparison lattice,
\(\sum_{r\in\mathbb Z\setminus\{0,1\}}1/[r(r-1)]=2\).
The correction that survives this cancellation can be read from the
first radial tail. Using \(\sum_{i\ge1}q_i^{-2}=1/2\) gives
\[
 2\sum_{i>K}(y_i^{-2}-q_i^{-2})
   =\frac{\Var_T U}{a^2}
      -2\sum_{i\le K}(y_i^{-2}-q_i^{-2}).
\]
To bound the three errors in this identity and the mixed sum, set
\(\beta=(\sqrt\kappa\,\ell^3)^{-1}\). On a disk of radius a fixed
multiple of \(\ell\), the inverse differs from the identity by
\(O(\kappa^{-1/2})\) and that difference vanishes to third order.
Cauchy's estimate on a smaller disk gives
\(|Y'(q)-1|\ll\beta |q|^2\). Extend the indices by
\(y_{1-j}=-y_j\). Integration between two levels now compares
their separation with its lattice value with relative error
\(O(\beta(q_i^2+q_j^2))\). In particular the positive radial head in
the last display contributes \(O(\beta K)\). The difference of the two
dimensionless mixed heads is bounded by
\[
 C\beta\!\!\sum_{\substack{1-K\le j\le K\\j\ne k,k+1}}
 \frac{j^2+(k+1)^2+1}{|j-k|\,|j-k-1|}
 \ll\beta(K+(k+1)^2).
\]
For the terms of positive degree in each reflected tail use
\(y_K\asymp K\), \(y_{k+1}=O(k+1)\), and
\(\sum_{j>K}y_j^{-2}=O(K^{-1})\). The last bound follows by
subtracting the counted head from the exact second trace as above,
using \(\sum_{j>K}q_j^{-2}=O(K^{-1})\) and \(\Var_T U/a^2=O(\ell^{-2})\).
The fourth-power tail is consequently \(O(K^{-3})\). Summing the even
degrees starting at two in \eqref{eq:exterior-tail} costs
\(O((k+1)^2 K^{-3})\) for each tail, since \(y_{k+1}/y_K<1/2\)
eventually. Finally \(a g_k=\pi(1+O(\beta(k+1)^2))\).
The only surviving term at scale \(\ell^{-2}\) for \(k+1=o(\sqrt\ell)\)
is \(\pi^2\Var_T U/a^2\). Substitution of
\(\kappa\Var_T U=1+O(\kappa^{-1/2})\) gives \eqref{eq:exterior-signed}.

Thus \(\ell^2(g_k^2M_k-2)\to\pi^2\) and \(D_k'<0\) uniformly on any
family \(k+1=o(\sqrt\ell)\) at sufficiently large \(T\). Section~\ref{sec:transport}
recovers the same leading drift on a wider large-time window by a phase estimate.
On a common large-time interval \([t,T]\), if \(t\to\infty\) and
\(\sup_{t\le s\le T}(k+1)^2/\ell_s\to0\), integration gives
\[
 D_k(t)-D_k(T)
     =(\pi^2+o(1))\int_t^T\frac{\dd s}{\ell_s^2}.
\]
This integration uses localization at each time in the interval. To obtain a time remainder from a single anchor we return to the field itself.

\subsection{The same data under backward evolution}
For \(s>0\) and \(P_m(z)=\prod_{i\le m}(1-b_i z^2)\), Gaussian convolution gives the exact earlier field
\begin{equation}
 \frac{H(T-s,z)}{Z(T)}
   =\E_T[e^{-sU^2}\cos(zU)]
   =\lim_{m\to\infty}\sum_{j=0}^m\frac{s^j}{j!}\partial_z^{2j}P_m(z).
 \label{eq:whole}
\end{equation}
Here \(b_i=x_i(T)^{-2}\) as in Section~\ref{sec:comparison}. Indeed \(|P_m(w)|\le H(T,\ii|w|)/Z(T)\); its log growth is \(O_T(1+|w|\log(2+|w|))\). Dominated convergence against the translated Gaussian therefore holds locally uniformly. The mass loss of the single frozen cosine with the exact initial curvature has the explicit size
\[
 \log\frac{Z(T-s)}{Z(T)e^{-s\E_T U^2}}
   =\int_0^s(s-v)\Var_{T-v}(U^2)\dd v .
\]
This follows by differentiating \(\log Z\) twice. Since \(\Var_T(U^2)\sim4L^2/\kappa\), the loss tends to one in logarithmic units at \(s=2\sqrt{T/L}\) \cite[\S4]{Founding07}. Equation~\eqref{eq:whole} gives a direct comparison for any backward estimate built from the same anchor data. Using a finite set of factors in that comparison requires a bound for the omitted product.

\section{Gap curvature from one upper time}\label{sec:anchored-curvature}

The exterior at \(T\) gives the exact initial slope in \eqref{eq:exterior-gap}.
To control the error of extrapolating this slope, tilt the complete law at
\(T\) by \(e^{-sU^2}\), as in \eqref{eq:whole}. Throughout, the saddle
parameters \(L,\kappa,S\) of \eqref{eq:saddle} are evaluated at \(T\).
Set
\[
 \lambda=L\sqrt\kappa,\quad r=\tfrac12\sqrt{\log\kappa},\quad
 \gamma=2L/\sqrt\kappa,\quad b=\gamma s,\quad
 D_*=\pi^2/(4L^2),\quad \eta=\kappa^{-3/4} .
\]
Here \(s\) is elapsed backward physical time, and \(y=Lx\) is the scaled
spatial argument. A joint complex estimate in \(b\) and \(y\) controls the
difference of the two inverse positions before time differentiation
\cite{EZCurvature}.

\Needspace{8\baselineskip}
\begin{theorem}[Signed remainder on a backward interval]\label{thm:anchored-curvature}
For all sufficiently large \(T\), the marked zeros for each \(k\ge1\) with
\(q_{k+1}\le\lambda r/8\) stay simple, adjacent and separated from the rest
of the divisor throughout \(0\le s\le s_*:=\gamma^{-1}\). Uniformly over
these pairs and times, with primes denoting derivatives in physical time,
\begin{align}
 \frac{D_k'(T-s)}{D_*}
   &=-\frac4\kappa\{1+b/\lambda+6\alpha b+O(\eta)\}<0,
      \qquad \alpha=\frac{2S}{3\kappa^{3/2}},\label{eq:anchored-slope}\\
 \frac{D_k''(T-s)}{D_*}
   &=\frac8{\kappa^2}+\frac{32LS}{\kappa^3}+O(L\kappa^{-9/4})
     \sim\frac{32L}{\kappa^2}.\label{eq:anchored-curvature}
\end{align}
In particular, the exact anchor slope \(E_k=D_k'(T)=2-g_k(T)^2M_k(T)\)
gives the signed estimate
\begin{equation}\label{eq:anchored-linear}
 \frac{D_k(T-s)-D_k(T)+sE_k}{D_*}
   =s^2\left\{\frac4{\kappa^2}+\frac{16LS}{\kappa^3}
                  +O(L\kappa^{-9/4})\right\} .
\end{equation}
It is strictly positive for every \(0<s\le s_*\).
For each Taylor degree \(m\ge2\) the bound with exact coefficients is
\begin{equation}\label{eq:anchored-taylor}
 \frac{\bigl|D_k(T-s)-\sum_{j=0}^m(-s)^j D_k^{(j)}(T)/j!\bigr|}{D_*}
 \le \frac{C\eta}{\lambda}\,\frac{(|b|/2)^{m+1}}{1-|b|/2}
 \qquad (0\le s\le s_*).
\end{equation}
The constant \(C\) is independent of the pair, Taylor degree and partial time.
\end{theorem}

\begin{proof}
\emph{The joint transform.}
Write \(W=\sqrt\kappa(U-L)\) and \(\zeta=y/\lambda\). Removing the scalar
\(e^{-sL^2}\) leaves the tilt \(e^{-bW-bW^2/(2\lambda)}\).
Use the unnormalized density \(F\) of Lemma~\ref{lem:saddle}, extended by
zero for \(w<-\lambda\), and the Gaussian model
\(\varphi(w)=e^{-w^2/2}/\sqrt{2\pi}\). Define
\[
 a_\pm=-b\pm\ii\zeta,\qquad v_b=A_b^{-1},\qquad A_b=1+b/\lambda,
\]
\[
 J_\pm=\int_{\mathbb R} F(w)e^{a_\pm w-bw^2/(2\lambda)}\dd w,
 \qquad I_\pm=v_b^{1/2}e^{v_ba_\pm^2/2},\qquad m_\pm=v_b a_\pm .
\]
Up to the removed scalar and a common nonzero anchor constant, the field is
\((e^{\ii y}J_++e^{-\ii y}J_-)/2\).
The integral with \(\varphi\) in place of \(F\) is \(I_\pm\), with its square
root continued from \(v_b=1\). Three derivatives in \(a_\pm\) show that
the Gaussian integral with the cubic correction in Lemma~\ref{lem:saddle}
is \(I_\pm\{1-\alpha(m_\pm^3+3v_bm_\pm)\}\).

For \(|b|\le3,\ |\zeta|\le r\), complete the real square in the absolute
central error integral on \(|w|\le w_0=\kappa^{1/12}\). Its tilted polynomial
moment is \(O(1+(r+3)^6)\). Division by \(|I_\pm|\) leaves a bounded
prefactor and the exponential of
\[
 \frac{(\Real a_\pm)^2}{2\Real A_b}-\frac12\Real(v_ba_\pm^2)
 =\frac12(\Imag a_\pm)^2+O((r+3)^2/\lambda).
\]
The resulting relative central error is \(O(\eta)\), since the density
error supplies \(\kappa^{-1}\) and the exponential costs at most
\(\kappa^{1/8+o(1)}\). The complete source on \(w\ge-\sqrt\kappa/4\)
outside the central interval is bounded by \(Ce^{-cw^2}\). The tilt
costs at most \(e^{(r+3)|w|+3w^2/(2\lambda)}\), leaving a fixed Gaussian
exponent beyond \(w_0\). The Gaussian model tails obey the same estimate,
including beyond the source's left endpoint. On
\(-\lambda\le w<-\sqrt\kappa/4\) the source is \(O(e^{-c\kappa})\) and
the integrated tilt costs at most
\(\lambda e^{(r+3)\lambda+3\lambda/2}=e^{o(\kappa)}\).
Thus all tail errors are smaller than any inverse power of \(\kappa\)
even after division by \(|I_\pm|\).

It follows that \(J_\pm/I_\pm=1-\alpha(m_\pm^3+3v_bm_\pm)+O(\eta)\). The deviation
from one is \(o(1)\) and its square is \(o(\eta)\), so
\begin{equation}\label{eq:anchored-log}
 \log J_\pm=\tfrac12\log v_b+\tfrac12 v_ba_\pm^2
                  -\alpha(m_\pm^3+3v_bm_\pm)+O(\eta).
\end{equation}
The integrals are holomorphic jointly by the same dominating bounds on
the complex disks. The remainder in \eqref{eq:anchored-log} is therefore
jointly holomorphic in the interiors and uniformly bounded by \(C\eta\).

\emph{The counted phase and its inverse.}
The field factors as a nonvanishing amplitude times \(\cos\Theta_b(y)\), where
\begin{equation}\label{eq:anchored-phase}
 \Theta_b(y)=y+\frac{\log J_+-\log J_-}{2\ii}
 =v_b y-\alpha\{(3v_b^3 b^2+3v_b^2)\zeta-v_b^3\zeta^3\}+O(\eta).
\end{equation}
The Gaussian part uses \(1-v_b b/\lambda=v_b\). The phase is odd in
\(y\), by exchange of \(J_+\) and \(J_-\), and real for real \(b,y\).
In the coordinate \(p=v_by\), the correction to \(p\) has size
\(O(\alpha(1+r^3)+\eta)\) on \(|p|\le3\lambda r/4\).
Rouché on \(|p|=\lambda r/2\) gives a unique simple inverse
\(y=Y_b(q)\) for \(|q|\le\lambda r/4\), jointly holomorphic with room at
the \(q\)-boundary. Cauchy's estimate on the larger \(p\)-disk gives
slope \(1+o(1)\) in \(p\) on the real segment through these inverses.
Thus \(\Theta_b(0)=0\) and the factorization count every positive zero
through \(Y_b(\lambda r/4)/L\) for real \(0\le b\le1\). In particular
the marked endpoints are \(Y_b(q_k)/L,Y_b(q_{k+1})/L\).

The inverse equation first gives \(p-q=O(\alpha(1+r^3)+\eta)\).
Replacing \(p\) by \(q\) in the cubic correction costs at most
\(O(\alpha(1+r^2)|p-q|/\lambda)=o(\eta)\). Replacing \(v_b\) by one
in that correction costs \(O(\alpha(1+r^3)/\lambda)=o(\eta)\). Hence
\begin{equation}\label{eq:anchored-inverse}
 Y_b(q)=q/v_b+\alpha\{(3b^2+3)q/\lambda-(q/\lambda)^3\}+R_Y(b,q),
 \qquad |R_Y|\le C\eta .
\end{equation}
Cauchy in \(q\) gives \(|\partial_q R_Y|\le C\eta/(\lambda r)\) on
\([0,\lambda r/8]\), uniformly in the time disk. Subtract
\eqref{eq:anchored-inverse} at the two adjacent levels before squaring.
In units of \(D_*\) the resulting holomorphic squared gap is
\begin{align}
 \mathcal D_k(b)&=\left(\frac{Y_b(q_{k+1})-Y_b(q_k)}{\pi}\right)^2
                   =P_k(b)+R_k(b),\label{eq:anchored-model}\\
 P_k(b)&=(1+b/\lambda)^2+\frac{2\alpha}{\lambda}
      \left(3b^2+3-\frac{q_{k+1}^3-q_k^3}{\pi\lambda^2}\right),
      \qquad |R_k(b)|\le C\eta/\lambda\quad (|b|\le3).\notag
\end{align}
The cross correction from replacing \(1/v_b\) by one in the order-\(\alpha\)
term after squaring costs \(O(\alpha(1+r^2)/\lambda^2)\). The square
of the gap correction costs \(O(\alpha^2(1+r^2)^2/\lambda^2)\).
Both fit the remainder along with the adjacent difference of \(R_Y\).

\emph{Physical time and clearance.}
For real \(s\) the squared gap satisfies \(D_k(T-s)=D_*\mathcal D_k(b)\).
The first two \(b\)-derivatives of \(R_k\) are \(O(\eta/\lambda)\) on
\(0\le b\le1\) by Cauchy. Using \(\mathrm d/\mathrm dt=-\gamma\,\mathrm d/\mathrm db\)
gives \eqref{eq:anchored-slope} and \eqref{eq:anchored-curvature}.
The two curvature terms come from the Gaussian part and the cubic source
correction, respectively. Their sum also equals the relative curvature
at \(T\) of the instantaneous lattice value \(\pi^2/(4L_t^2)\), by
\(L_t'=2L_t/\kappa_t\) and \eqref{eq:saddle}. Integration against \((s-u)\dd u\) for
\(0\le u\le s\) proves \eqref{eq:anchored-linear}. Subtracting the degree-\(m\)
Taylor polynomial of \(\mathcal D_k\) leaves only the Taylor tail of \(R_k\)
when \(m\ge2\). Its radius-two Cauchy bound gives \eqref{eq:anchored-taylor}.

The slope and \eqref{eq:anchored-model} give \(D_k(T-s)\ge D_k(T)>D_*/2\).
The outer level \(q_{k+2}\) lies within \(\lambda r/4\) eventually.
The next zero on the left is counted there as well, with the reflected
first zero used for \(k=1\). The uniform \(o(1)\) displacement from \(q/v_b\)
in \eqref{eq:anchored-inverse} gives adjacent spacings
\((\pi/L)(1+o(1))\). The field is all real at these large physical times
\cite{Polymath2019}, so a circle of radius \(g_k\) about the real midpoint
contains exactly the pair and stays clear of the other zeros.
\end{proof}

\subsection{The anchor coefficients and their precision}
The complete hierarchy identifies the derivatives in
\eqref{eq:anchored-taylor} on the separated chart. With the pair index
suppressed, combine \eqref{eq:exterior-gap} with the diagonal equation
in \eqref{eq:parent} to give
\begin{equation}\label{eq:anchored-polynomial}
 D''=-8M+8DM^2+48D S_{2,2}+64D^2 S_{3,3} .
\end{equation}
The unreduced integer split recurrence raises the largest total pole
order of a row by at most two per differentiation. Applied to \(D'=2-4DM\),
it uses rows of order at most \(2m\) for the Taylor jet of degree \(m\).
The diagonal reduction in \eqref{eq:anchored-polynomial} uses an order-six
row \cite[\S3]{AGC}. The joint disk bounds the resulting scalar gap. An
enclosure for an individual evolving row requires its own analytic estimate.

For the finite anchor calculation use \(y_i=Lx_i(T)\) and take the head
through \(q_K\le\lambda r/4\) with \(K\) maximal. The marked endpoints have
an upper bound \(A\) and the unknown positive roots a lower bound \(R\),
with \(R\asymp K\) and \(A/R\le2/3\) eventually by \eqref{eq:anchored-inverse}.
Apply \eqref{eq:exterior-tail} in these coordinates, with
\(W_j^{(y)}=\sum_{i>K} y_i^{-2j}\) and moments \(\nu_j=\E_T(U/L)^j\).
For any fixed row of order \(h=p+q\), retaining degrees \(d<N\) uses
moments through order at most \(h+N-1\). Each known exterior occurrence
in that row is evaluated with positive lower distances \(l_+,l_-\).
For position errors \(e_i,e_+,e_-\) its error is at most
\[
 \frac{p(e_i+e_+)}{l_+^{p+1}l_-^q}
   +\frac{q(e_i+e_-)}{l_+^p l_-^{q+1}},
\]
with terms of zero numerator omitted. The same endpoint enclosures enter
each row and the gap. Widths from the radial-power subtractions and
the target coefficients in the retained tail must be propagated as well.

Here is a sufficient acquisition budget for the slope and curvature
\cite[finite-precision supplement]{EZCurvature}. In the scaled variables
\(D_y=L^2D\), \(M_y=M/L^2\) and \(S_{p,q}^{(y)}=S_{p,q}/L^{p+q}\), use
\(\sigma=5/6\) in the bound following \eqref{eq:exterior-tail} and choose
\(N=O(\log\lambda)\) with \((4/5)^N\le\lambda^{-6}\). The omitted tail
at orders \(h=2,4,6\) is \(O(\lambda^{-6})\), since
\(W_1^{(y)}\le\nu_2/2=O(1)\) and \(R^{2-h}\le1\) eventually. The retained
degrees need moments through order \(N+5\).

On the unit disk the truncated radial cosine polynomial lies within
\(3/4\) of one, with strict clearance for sufficiently accurate input
enclosures, since the corresponding full positive coefficient sum tends
to \(\cosh 1-1\). This limit follows by the weighted saddle bounds. If
\(\delta\) is the sum of absolute errors in its normalized coefficients,
Cauchy's estimate for the perturbed logarithm gives an error
\(O(j\delta)\) in the full \(j\)-th radial power. A common position radius
\(e_{\rm pos}\) for the scaled head gives an additional \(O(je_{\rm pos})\)
in the head subtracted from it, using \(y_i-e_{\rm pos}\ge i\) and summing
the derivative bounds for \(y_i^{-2j}\). The retained coefficients satisfy
\(|h_d|\le\binom{h+d-1}{d}A^d\), so the summed error from these radial
subtractions is at most
\begin{equation}\label{eq:anchored-input-width}
 O\big((\delta+e_{\rm pos})N^7\max(1,A)^N\big).
\end{equation}
With the exact radial powers held fixed, the endpoint errors in the tail
cost \(O(e_{\rm pos})\) by differentiating the convergent reflected sums
under the radial clearance. The counted exterior head has the same bound
from its separated spacings and summable reciprocal products.

Thus \(\delta,e_{\rm pos}\ll\lambda^{-6}/(N^7\max(1,A)^N)\) suffice,
which calls for \(O(\log^2\lambda)\) absolute bits in these scaled data.
The inputs of the polynomial for \(D''/L^2\) in
\eqref{eq:anchored-polynomial} are bounded in the scaled coordinates.
Errors \(O(\lambda^{-6})\) persist at that order in the polynomial and in
\(E_k=2-(Lg_k)^2M_y\). The curvature signal at the anchor is
\(D''/L^2\sim8\pi^2L/\lambda^4\), which dominates this error.

The full odd phase also gives the head to the required precision. In
\(y\), its linear coefficient is \(a/L\). The remaining odd part is bounded
by \(M_{\rm ph}=O(\alpha(1+r^3)+\eta)\) on \(|y|\le\lambda r\), by
\eqref{eq:anchored-phase} at \(b=0\) and Cauchy for the linear term.
On \(|y|\le X:=\lambda r/2\), its Taylor tail past degree \(P\) is
\(O(M_{\rm ph}2^{-P})\). Take \(P=O(\log^2\lambda)\) to make the tail
smaller than \(e_{\rm pos}\). The coefficients are obtained from the log
of the truncated moment generating polynomial for \(U/L\). On the disk
of radius \(1/4\), the positive coefficient sum excluding the constant
tends to \(e^{1/4}-1<1/2\) when all moments are included. Thus an error
\(\delta_{\rm ph}\) in the sum of coefficient radii weighted at \(1/4\)
costs \(O(4^j\delta_{\rm ph})\) in the \(j\)-th log coefficient. It suffices
that \(\delta_{\rm ph}\ll e_{\rm pos}(4X)^{-P}\). The true phase slope
is bounded below and the counted head lies strictly inside \([0,X]\), so
the resulting \(O(e_{\rm pos})\) phase error bounds the positional error.
This gives a sufficient moment order \(O(\log^2\lambda)\) and
\(O(\log^3\lambda)\) absolute bits in the stated weighted coefficient norm,
including acquisition of the head. These are asymptotic tolerances.
An application at a specified \(T\) requires checking the inequalities
with the actual input enclosures.

For a finite Taylor calculation write \(d_j\) for the absolute error
in \((-1)^j D_k^{(j)}(T)/(D_*\gamma^j j!)\), allowing for scale uncertainty
and rounding. Add \(\sum_{j=0}^m d_j|b|^j\) to \eqref{eq:anchored-taylor}.
For example, \(\sum_{j=2}^m d_j=O(\eta/\lambda)\) adds a normalized
error \(O((\eta/\lambda)b^2)\) for \(|b|\le1\).
In \eqref{eq:anchored-linear} inexact anchor values add physical errors
\(|\delta D_k(T)|+s|\delta E_k|\). Resolving its signed quadratic increment
at a specified \(b>0\) requires \(d_0+d_1|b|\) small compared with
\((\alpha/\lambda)b^2\). Exact shared anchor values yield the \(s^2\)
vanishing down to \(s=0\). The finite coefficient budget is for the
anchor data, while \eqref{eq:anchored-model} supplies the complete analytic
time allowance from the source.

\subsection{Count, scope and comparison}
The proof counts \(K=\lfloor\lambda r/(4\pi)+1/2\rfloor\) positive zeros
through \(Y_b(\lambda r/4)/L\). The retained pairs number
\(\lfloor\lambda r/(8\pi)+1/2\rfloor-1\), of order
\(\sqrt{T}(\log T)^2\), on a physical interval of length
\(s_*\sim2\sqrt{2T/\log T}=o(T)\). The onset in \(T\) is unevaluated.
The current estimate of Section~\ref{sec:transport} treats a fixed multiple
of the saddle window using bounds along the interval. The present estimate
uses one upper law and a time disk on a growing window. A separate spatial
estimate treats sufficiently high labels on a fixed slab in \((0,1/2]\),
with its own direct two-time comparison \cite[\S5 and references there]{EZCurvature}.
A backward path from the present upper anchor to a fixed lower time
requires different bounds.

Direct evolution by \eqref{eq:whole} uses the same complete law. Integrating
\(\Theta_b'\) between the two counted roots gives \(\pi\), so the direct
phase calculation with the same spatial Cauchy estimate and time disk
recovers \eqref{eq:anchored-model} and the same curvature and Taylor
bounds. The shared analytic error provides the gain over separate endpoint
intervals, whose widths can persist as \(s\) decreases to zero.
The anchor estimate identifies these bounds with the derivatives from
the exterior hierarchy.

The sign uses the quantitative tilted theta expansion. For example, a
stationary cosine has zero gap curvature and its atomic law lacks the cubic
saddle correction. Arithmetic and profile modifications change the complete
source, and changed-observation controls fail the literal heat-cosine tilt.
The tail identity \eqref{eq:exterior-tail} requires a complete all-real
divisor with integer multiplicities. Controls with nonreal zeros or relaxed
divisor weights fail this premise. Finally, a backward collision of the
atomic field \(e^t\cos x+c\) destroys the separated chart
\cite[\S5 and references there]{EZCurvature}.

\section{Gap motion from the phase current}\label{sec:current}
The upper-anchor curvature theorem controls a time interval by a single complex tilt. A different identity computes the instantaneous gap velocity from the spatial curvature of a real current. It also determines where this phase coordinate can be used.

\subsection{The current and its domain}
On a real interval with \(B\ne0\), take a continuous argument \(\Theta\) of \(B\) and write \(B=e^{p+\ii\Theta}\). At the origin \(\Theta=0\). Define the current and its normalization by
\begin{equation}
 J(t,x)=\Imag(\overline{B(t,x)}B_x(t,x))=|B(t,x)|^2\Theta_x(t,x),
 \qquad j(t,x)=J(t,x)/(Z(t)^2a_t).
 \label{eq:current}
\end{equation}
The complete source determines the current even where the real part \(H\) vanishes. A positive current on \([0,X]\) gives a nonvanishing \(B\) and strictly increasing phase there. To include a labelled edge one must also reach its upper level \(q_{k+1}\).

For this particular lift one can identify a whole time range where the
coordinate fails between two positive real zeros. Put
\(t_{\rm m}=-\Phi''(0)/(2\Phi(0))>5/2\), where the last inequality
uses \(\Phi''(0)<-10\) and \(0<\Phi(0)<2\) \cite{EZPhaseScope}.
By \href{https://doi.org/10.1007/BF02075457}{Csordas and Varga, Theorem 2.1},
\(v\mapsto\log\Phi(\sqrt v)\) is strictly concave and its right derivative at
zero is \(-t_{\rm m}\). Hence \(f_t(u)=e^{tu^2}\Phi(u)\) is strictly
decreasing for \(u>0\) when \(t\le t_{\rm m}\). For \(x>0\), put \(h=\pi/x\).
Pairing half-periods gives
\[
 \Imag B(t,x)=\sum_{j\ge0}\int_0^h
 \{f_t(2jh+u)-f_t((2j+1)h+u)\}\sin(xu)\dd u>0 .
\]
The upper-half-plane argument consequently takes the value \(\pi/2\) at
both real zeros. It is nonconstant between them by analyticity of \(H\)
and \(H(t,0)>0\). Thus \(J\) takes both signs between the endpoints.
At any fixed time the same obstruction applies to sufficiently high positive
zeros, since integration by parts gives
\(\Imag B(t,x)=\Phi(0)/x+O_t(x^{-2})\).
The argument concerns the half-line lift. The exterior calculation in
Section~\ref{sec:exterior} uses its own zero counts and separation bounds.

\begin{theorem}[Velocity, interaction and joint forcing]\label{thm:velocity}
On a common increasing phase chart define \(v=(\log j)_x\) and
\(\mathcal A=(v(t,x_+)-v(t,x_-))/g\). Then
\begin{align}
 \dot x_\pm&=v(t,x_\pm),&
 (\log g)^{\displaystyle\cdot}&=\mathcal A
   =\frac1g\int_{x_-}^{x_+}(\log j)_{xx}\dd x, \label{eq:velocity}\\
 M&=\frac1{2D}-\frac{\mathcal A}{2},&
 \mathcal F&=\frac1{2D^2}+\frac12(\dot{\mathcal A}+\mathcal A^2).
 \label{eq:forcing}
\end{align}
The same \(v\) advances every edge in the chart. In the parent's normalization \(V=1-2DM=D\mathcal A\) and \(K=2D^2\mathcal F=1+D^2\ddot g/g\).
\end{theorem}
\begin{proof}
The imaginary part of \(B_t/B=-B_{xx}/B\) gives
\(\Theta_t=-\Theta_{xx}-2p_x\Theta_x\). Differentiating
\(\Theta(t,x_\pm(t))=q_{k},q_{k+1}\), respectively, yields
\(\dot x=2p_x+\Theta_{xx}/\Theta_x=(\log J)_x\). Its difference proves \eqref{eq:velocity}. Thus \(\dot D=2D\mathcal A\); substituting in both equations of \eqref{eq:parent} gives \eqref{eq:forcing}, including the cancellation of the two terms linear in \(\mathcal A/D\).
\end{proof}
For the forcing there is also an endpoint expression free of separate exterior sums. With \(b_v=v_t+v v_x\), one has
\begin{equation}
 K-1=\frac{D^2}{g}\{b_v(t,x_+)-b_v(t,x_-)\}.
 \label{eq:acceleration}
\end{equation}
Differentiate the two velocity evaluations along their actual paths. The time derivative of \(v\) is \(\partial_x(J_t/J)\). The unnormalized double-source formula below computes \(J_t\) by inserting \(u^2+u'^2\) in the same integral. A signed bound for this difference could therefore feed the parent's joint forcing barrier directly. A bound on \(\mathcal A\) alone does not bound its time derivative.

\subsection{A positive source for spatial curvature}
Take independent copies \(U,U'\) under \(\E_t\), and use the probability weight \((U+U')/(2a)\) on their joint law. Write \(\E_*\) for this weighted expectation and \(\Delta=U-U'\). Symmetrizing \(\overline B B_x\) proves
\begin{align}
 j(t,x)&=\E_*\cos(x\Delta),&
 w&=\E_*\Delta^2
 =\frac{\E U^3-a\E U^2}{a}=\frac{\dot a}{a},&
 m_4^*&=\E_*\Delta^4 . \label{eq:difference}
\end{align}
The transform \(j\) is even and entire, and \(j(0)=1\). The small frequency in this integral is the difference of the two source frequencies. In particular the quadratic term of \(\log j\) is exactly \(-wx^2/2\).

\begin{proposition}[A finite two-sided gap estimate]\label{prop:finite}
At each time in a common counted chart with endpoints in \([0,X]\), put
\[
 b_* =1-X^2 w/2,\qquad
 l=w-X^2m_4^*/2,\qquad u=w/b_*+X^2w^2/b_*^2 .
\]
If \(b_*>0,\ l\ge0\), then \(l\le-\mathcal A\le u\).
Consequently for \(t\le T\) with these inputs throughout,
\begin{equation}
 D(T)\exp\!\left(2\int_t^T l(r)\dd r\right)
 \ \le D(t)\le\
 D(T)\exp\!\left(2\int_t^T u(r)\dd r\right).
 \label{eq:slab}
\end{equation}
\end{proposition}
\begin{proof}
On \(0\le x\le X\) one has \(b_*\le j\le1\), \(|j'|\le xw\), and
\(-w\le j''\le-w+x^2m_4^*/2\), using \(\cos y\ge1-y^2/2\).
Insert these signed bounds into \(j''/j-(j'/j)^2\), then integrate \eqref{eq:velocity}. The condition \(l\ge0\) permits division of the negative upper bound by \(j\le1\).
\end{proof}

\subsection{The higher exterior moments from one deflated field}
At a fixed counted pair put \(Q(z)=(z-c)^2-D\) and
\(R=\partial_z\log(H/Q)\), analytic near the target segment in a cleared chart. Set \(r^2=D+\xi\), with \(\xi\) small enough that the probes \(c\pm r\) stay in that chart and \(|\xi|<\inf_{\rho\ {\rm ext}}|Q_\rho|\). The parent's complete rational traces \cite[\S3]{AGC} then give
\begin{align}
 \sum_{n\ge0}S_{n+1}\xi^n
 &=-\frac{R(c+r)-R(c-r)}{2r}, \label{eq:Sgen}\\
 \sum_{n\ge0}T_{n+2}\xi^n
 &=-\frac{R(c+r)+R(c-r)-R(c+\sqrt D)-R(c-\sqrt D)}{2\xi},\label{eq:Tgen}
\end{align}
where \(T_n=\sum_{\rho\ {\rm ext}}(\rho-c)Q_\rho^{-n}\) for \(n\ge2\). Quotients take their removable values where necessary. In fact an exterior pole at \(c+d\) contributes \((d^2-r^2)^{-1}\) to \eqref{eq:Sgen}. The symmetric sum in \eqref{eq:Tgen} before subtraction contributes \(-2d/(d^2-r^2)\); its subtraction at \(\xi=0\) removes the first-order ambiguity. The resulting series converge absolutely by the parent's reciprocal-square tail bound. These series generate the indicated diagonal and midpoint-weighted traces at a fixed pair. Section~\ref{sec:exterior} uses both endpoint exponents independently. Controlling its coefficients at lower time requires a positive radius and a bound on the same deflated field.

\subsection{A comparison from the source at each time}\label{sec:transport}
The same centered source estimate that locates the zeros in Section~\ref{sec:anchors} gives a comparison for their motion throughout a high-time interval. This calculation uses the phase velocity of Section~\ref{sec:current}.

\begin{theorem}[Signed backward transfer]\label{thm:transfer}
Fix \(C>0\). There is a time \(t_C\) such that, for \(T\ge t\ge t_C\) and a marked edge satisfying \(q_{k+1}\le C\inf_{t\le r\le T}\ell_r\), the same two positive phase levels are simple throughout. Uniformly on that family,
\begin{align}
 (\log g)^{\displaystyle\cdot}&=-2/\kappa_r+O_C(\kappa_r^{-2})<0,
 &M&=\frac1{2D}+\frac1{\kappa_r}+O_C(\kappa_r^{-2}),\label{eq:asymM}\\
 D(t)&=D(T)(L_T/L_t)^2
           \exp\!\left(O_C\!\left(\int_t^T\kappa_r^{-2}\dd r\right)\right).
 \label{eq:asymD}
\end{align}
For any \(X_r\) containing the endpoints in \([0,X_r]\) and satisfying \(X_r\le C'\sqrt{\kappa_r}\), there is also a mean-normalized bound
\begin{equation}
 \left|\log\frac{a_t g(t)}{a_Tg(T)}\right|
 \le C''\int_t^T X_r^2\kappa_r^{-3}\dd r. \label{eq:meanD}
\end{equation}
The constants depend only on the fixed window. The asymptotic theorem gives existence of \(t_C\); Section~\ref{sec:finite} treats a specific interval independently.
\end{theorem}
\begin{proof}
At a physical spatial argument \(x=\sqrt\kappa\,\zeta\), the analytic continuation of the normalized current \eqref{eq:current} is
\[
 j(t,\sqrt\kappa\,\zeta)
 =\exp\{K_0(\ell\zeta)+K_0(-\ell\zeta)\}\,[1-\Psi'(\ell\zeta)].
\]
This follows on the real axis from the amplitude and phase in \eqref{eq:odd} and then by continuation. The odd cubic in \eqref{eq:complex-saddle} cancels from the exponential. Cauchy's estimate on a slightly larger disk gives \(\Psi'(\ell\zeta)=O_C(1/(\ell\sqrt\kappa))\). Hence on every fixed complex disk, for large enough \(t\),
\begin{equation}
 \log j(t,\sqrt\kappa\,\zeta)=-\zeta^2+O_C(\kappa^{-1}).\label{eq:logj}
\end{equation}
The branch is real at zero. Two spatial derivatives give
\((\log j)_{xx}=-2/\kappa+O_C(\kappa^{-2})\). The inverse theorem places the full real edge in a fixed multiple of \([0,\sqrt\kappa]\) and identifies its index at each time; uniqueness and simplicity give the same moving endpoints. Average the curvature in \eqref{eq:velocity} and use \eqref{eq:forcing}. From \eqref{eq:saddle}, \(\dot L=2L/\kappa\). Integrating the estimate for \(2\mathcal A\) gives \eqref{eq:asymD}. The integrated remainder is \(O_C(1/[t(\log t)^2])\) even if the upper time tends to infinity.

Four spatial derivatives of \eqref{eq:logj} give \((\log j)_{xxxx}=O_C(\kappa^{-3})\). Evenness and \eqref{eq:difference} give \((\log j)_{xxx}(0)=0\) and \((\log j)_{xx}(0)=-\dot a/a\). Two real integrations therefore give \(|(\log j)_{xx}+\dot a/a|\ll_C x^2/\kappa^3\). Average over the edge and integrate in time, proving \eqref{eq:meanD}. For fixed \(k\), \(X_r=O_k(1/a_r)\) is allowed and the last integral is \(O_k(1/[t^2(\log t)^5])\).
\end{proof}
These are estimates on the complete interaction and the signed drift of the actual edge. The individual source expansions also localize its endpoints throughout the same high-time domain. The estimate establishes direction of motion and its integrated size there. Extending it to an earlier region requires additional source information on the intervening path.

\paragraph{The correction after removing the mean spacing.}
Let \(h=\log j\) and denote spatial averaging over the marked interval by
angle brackets. Combining \eqref{eq:velocity} with \(h_{xx}(t,0)=-\dot a/a\)
and \(h_{xxx}(t,0)=0\) gives the exact identity
\begin{equation}
 \frac{\dd}{\dd t}\log(ag)=
  \left\langle\int_0^x(x-y)\,\partial_y^4h(t,y)\dd y\right\rangle,
 \qquad \partial_x^4h(t,0)=m_4^*-3w^2 .\label{eq:quartic}
\end{equation}
Freezing the fourth derivative at zero yields
\((m_4^*-3w^2)(x_-^2+x_-x_++x_+^2)/6\).
If \(|\partial_x^6 h|\le B_6\) on \([0,X]\), evenness gives an
absolute error at most \(B_6 X^4/24\). This retains the cancellation between
the fourth moment and the squared variance when one compares a change
of gap with the change of its natural mean spacing.

\section{A finite comparison for the first gap}\label{sec:finite}
\subsection{A uniform bound from finite data at one time}
Take the first positive edge, \(X=4\), \(s=1000-t\in[0,1]\), and \(b=671/500\). Write \(P=U-b,\ W=U^2-b^2\), and let \(\E_0\) here denote expectation at the upper anchor \(1000\). Tilting this one source gives the moments
\begin{equation}
 d_j(s)=\frac{\E_0[P^je^{-sW}]}{\E_0[e^{-sW}]},\quad d_0=1,\quad
 \mu=b+d_1,\quad \sigma^2=d_2-d_1^2. \label{eq:tilt}
\end{equation}
The numerator for each \(0\le j\le6\) is enclosed by a cubic polynomial with absolute remainder
\begin{equation}
 \left|\E_0[P^je^{-sW}]
    -\sum_{m=0}^3\frac{(-s)^m}{m!}\E_0[P^jW^m]\right|
 \le\frac{s^4e^{b^2}}{24}\sqrt{\E_0[P^{2j}W^8]} . \label{eq:tiltrest}
\end{equation}
Indeed \((-sW)_+\le b^2\), and Taylor's integral remainder followed by Cauchy--Schwarz gives \eqref{eq:tiltrest}. Since \(W=P(2b+P)\), all source integrals are anchor moments through degree 28. The two-copy law gives, for \(p=2,4\),
\begin{equation}
 m_p^*=\frac1\mu\sum_{i=0}^p(-1)^i\binom pi d_i
                    \{b\,d_{p-i}+d_{p-i+1}\}. \label{eq:certmom}
\end{equation}
This expands \(\E[U(U-U')^p]/\mu\), using the exchange symmetry for even \(p\).

All moments were evaluated by outward Arb arithmetic at 224 bits, using the original anchor integrator and its complete theta tail allowance. It integrates the first square mode on \([1,2]\), bounds the higher square modes there by geometric comparison, and bounds the two omitted half-line intervals by the saddle estimates in \cite[\S6]{Founding08}. Its absolute allowance is below \(2.353\,10^{-107}\) in each peak-scaled integral for functions bounded by \((u+2)^7e^{16(u+2)}\). This covers \(|P|^{28}\) and all lower powers, since \(21\log(u+2)\le (21/e)(u+2)\). It also covers the real trigonometric endpoint tests below.

Evaluate \eqref{eq:tiltrest} on 16 equal closed subintervals of \(s\). Each tilted mass has a positive lower bound. On every subinterval \(j(t,x)>0.99817\) throughout \([0,4]\). Also \(\left|\E e^{\ii x(U-\mu)}-1\right|\le x^2\sigma^2/2\); the checked inequality \(\mu X-\arcsin(X^2\sigma^2/2)>3\pi/2\) places both labelled zeros inside the interval. The centered transform remains in the corresponding disk around one, which fixes the argument from zero. These are whole-cell inequalities.

\Needspace{9\baselineskip}
\begin{proposition}[Finite signed descent]\label{prop:check}
For \(999\le t\le1000\) the first positive edge satisfies
\begin{equation}
 0.0002270493<-\dot g(t)/g(t)<0.0002300072.
\end{equation}
Its anchor value is \(D(1000)=1.3709993248674981982\ldots\).
Integrating the cell-specific bounds in Proposition~\ref{prop:finite} gives
\begin{equation}
 1.37162235<D(999)<1.37162972. \label{eq:number}
\end{equation}
\end{proposition}
\begin{proof}
Apply \eqref{eq:certmom} on each cell, check \(b_*>0,l>0\), and multiply the anchor \(D\) by the interval exponentials in \eqref{eq:slab}. The cell widths are exact dyadics. The program, precision, sign tests and outward endpoints are recorded alongside the manuscript \cite{FiniteRecord}.
\end{proof}
An equally informed endpoint calculation uses the same \(d_0,\ldots,d_6\) at \(s=1\). At real \(x\) evaluate
\[
 \sum_{j=0}^5\frac{x^jd_j}{j!}\cos(bx+j\pi/2)
 \quad\hbox{with error at most } |x|^6 d_6/720 .
\]
Signs on opposite sides and the common phase count enclose each root. This gives \(1.3716257900<D(999)<1.3716262696\), width below \(4.80\,10^{-7}\), versus \(7.36\,10^{-6}\) for transport. Direct quadrature of the complete upper-anchor field with the multiplier \(e^{b^2-U^2}\), using separate endpoint signs and the same count, gives \(D(999)=1.37162603031689388115504\ldots\). It shares source and tail code with the moment calculation. The path estimate establishes strict widening at every intermediate time. For the endpoint alone the finite Taylor comparison is tighter.

\subsection{Removing the change in mean spacing}
On a counted chart let \(h=\log j\) and let angle brackets denote spatial averaging along the edge. From \eqref{eq:difference} and evenness,
\begin{equation}
 \frac{\dd}{\dd t}\log(a g)
 =\left\langle\int_0^x(x-y)h^{(4)}(t,y)\dd y\right\rangle,\qquad
 h^{(4)}(t,0)=m_4^*-3w^2 . \label{eq:quartic}
\end{equation}
The superscript denotes spatial differentiation. Freezing the last derivative at zero gives the explicit edge term
\((m_4^*-3w^2)(x_-^2+x_-x_++x_+^2)/6\).
If \(\sup_{[0,X]}|h^{(6)}|\le B_6\), evenness and Taylor's theorem bound the error by \(B_6 X^4/24\). Thus a shared signed fourth cumulant and a controlled sixth derivative can replace separate absolute payments for \(j''\) and \((j')^2\). A later calculation with the same 28 upper-time moments carried this comparison to \(t=984\) for the first edge. Direct evaluation gave a tighter endpoint interval. Both methods certified widening along this already large-time path \cite{EZJoined}.

The material acceleration on this chart satisfies \eqref{eq:acceleration}. A bound on \(K\) together with initial \(D,V\) enters the parent's stopped-path barrier \cite[\S4]{AGC}. The divided differences \eqref{eq:Sgen}--\eqref{eq:Tgen} give a separate way of comparing that forcing with fixed-order exterior moments. Through a loss of the phase chart the parent's count-two contour \cite[\S\S2,6]{AGC} still requires a complete boundary floor and a paid error for its center and discriminant. Relabelling a new pair at the end of a step would invalidate the comparison.

A simple control tests these hypotheses. For \(H(t,x)=e^t\cos x+c\), \(c>1\), a positive half-line lift is \(B=e^{t+\ii x}+c\). Its current is \(e^{2t}+ce^t\cos x\). At \(ce^{-t}=1/2\), the first two roots are \(2\pi/3,4\pi/3\), and their physical gap velocity is \(2/\sqrt3>0\). The current is positive there, but \(m_4^*=w=1/3\); any \(X\ge4\pi/3\) fails \(w-X^2m_4^*/2>0\). At \(t=\log c,\ x=\pi\) the current vanishes as the two roots meet. The single-cosine control has constant positive current, \(w=0\), and stationary gaps. Both satisfy the representation identities. The signed theta estimate uses quantitative source concentration. Weighted-square and finite-prime controls change the normalized theta weights, reflected smooth controls lack the integer-square expansion, and altered observations change the cosine readout. The source and count assumptions for the additional controls are detailed in \cite{EZPhaseScope,EZMixedAnchor}. The finite computation is inside a previously all-real regime. Lower-time and all-edge coverage remain open.

\section{What sharper locations would give}\label{sec:frontier}
For the cubic approximation in \eqref{eq:shared-residual}, the first question is how much of its position error is shared between neighboring roots. The same source mean and odd cumulants enter both predictions. A useful estimate would preserve that dependence when taking the difference, and give a uniform error in units of the local spacing. Its comparison is the best bound obtained by differencing or differentiating the existing inverse remainder with the same source data. A separate way to gain height would be to quantify the onset of the inverse as the window constant grows. The complete source must then keep the centered transform away from zero on the chosen complex disk and preserve the count of its inverse.

At one specified all-real \(T\), the product formula supplies a second route to high labels. The finite-order inequalities \eqref{eq:factor-interval} give starting bounds on an individual zero. The question is how the required moment order and precision grow with \(k\) to put two consecutive zeros in intervals narrower than their local spacing. At \(T=1/2\), Section~\ref{sec:comparison} provides a direct comparison with a published counted approximation. Any information from that count used to initialize a product calculation must enter the comparison on both sides.

For a counted family at \(T\), Section~\ref{sec:exterior} makes the accuracy required of a location estimate visible in \(g_k^2 M_k-2\). The leading lattice terms cancel. The error in the complete expression, including zeros outside the located family, must be smaller than the surviving term to establish its sign. A backward estimate from this single time would then have to control
\[
 D_k(T-s)=D_k(T)-\int_0^s
       \bigl(2-g_k(T-u)^2S_{k,1,1}(T-u)\bigr)\dd u
\]
at every partial time, with a common count isolating the same pair. Repeated time differentiation of the mixed moments raises their orders. A usable time remainder must bound those higher orders and all unlocated zeros as the time interval lengthens. The Gaussian evolution of the same upper-time product in \eqref{eq:whole} supplies a direct comparison. A conclusion about all zeros would also require the central pair, the intervening heights and a count excluding nonreal zeros. For the bound in Section~\ref{sec:comparison}, the question is which zero-free region this information establishes at stronger parameters.

\paragraph{Acknowledgment.} Drafted for Edward Baker with substantial OpenAI Codex LLM assistance. The opening draft recorded gpt-6.2-astra at max effort. The present analytical revision used gpt-6.2-astra and gpt-6.2-sol at ultra effort. The cited project assessments are internal.
\begingroup
\small
\bibliographystyle{plain}
\bibliography{branch,references}
\endgroup
\end{document}
